%!TEX program = lualatex
%!TEX encoding = UTF-8 Unicode  
%!TEX root = chapter19.tex  
\documentclass[main.tex]{subfiles} 
\begin{document} 
	%----------------------------------------------------------------------------------------
	%	CHAPTER 2
	%----------------------------------------------------------------------------------------
	\cleardoublepage 
	\loadgeometry{marginlayout}  
	\pagestyle{scrheadings} % Print headers again  
	\KOMAoptions{
		headwidth=textwithmarginpar,
		headsepline=true,
		headsepline= 0.5pt,
		footwidth=textwithmarginpar,
	}   
	\onehalfspacing 
	\setcounter{chapter}{18}
	
%----------------------------------------------------------------------------------------
%	 
%----------------------------------------------------------------------------------------	
\chapter{Géométrie plane}

%----------------------------------------------------------------------------------------
%	 
%----------------------------------------------------------------------------------------



\section{Triangles égaux et semblables}

\begin{definition}[triangles égaux]
	Deux triangles sont égaux si leurs trois côtés et leur trois angles sont égaux deux à deux. 
\end{definition} 
%Les triangles $ABC$ et $IJK$ de la figure \ref{chapitre05:trianglesegaux} sont égaux. On a les 6 égalités suivantes :
%\begin{figure}[h]\centering
%	\begin{tikzpicture}[scale=0.9]
%		\tkzDefPoint(0,0){A}
%		\tkzLabelPoint[above left](A){$A$}
%		\tkzDefShiftPoint[A](-135:2.5){B}  
%		\tkzDefTriangle[two angles = 80 and 60](A,B)
%		\tkzGetPoint{C}
%		\tkzLabelPoint[left](B){$B$}
%		\tkzLabelPoint[right](C){$C$}
%		\tkzDrawPolygon(A,B,C)
%		\tkzDrawPoints(A,B,C)
%		\tkzDefShiftPoint[A](4,+3.0){J}  
%		\tkzDefShiftPoint[J](-70:2.5){I}
%		\tkzDefTriangle[two angles = -80  and -60](I,J)
%		\tkzGetPoint{K}
%		\tkzDrawPolygon(I,J,K)
%		\tkzDrawPoints(I,J,K) 
%		\tkzLabelPoint[right](I){$I$}
%		\tkzLabelPoint[left](J){$J$}
%		\tkzLabelPoint[right](K){$K$}
%		\tkzMarkAngle[mark=|,mkcolor=black, thick, size=0.6, fill=red!20, opacity=0.6](B,A,C)
%		\tkzMarkAngle[mark=||,mkcolor=black, thick,size=0.6, fill=blue!20, opacity=0.6](J,K,I)
%		\tkzMarkAngle[mark=||,mkcolor=black, thick,size=0.6, fill=blue!20, opacity=0.6](A,C,B)
%		\tkzMarkAngle[mark=|,mkcolor=black, thick,size=0.6, fill=red!20, opacity=0.6](K,I,J)
%		\tkzMarkAngle[mark=s,mkcolor=black, thick,size=0.6, fill=green!20, opacity=0.6](C,B,A)
%		\tkzMarkAngle[mark=s,mkcolor=black, thick,size=0.6, fill=green!20, opacity=0.6](I,J,K) 
%		\tkzDrawSegment[line width=2pt, color=red](B,C)
%		\tkzDrawSegment[line width=2pt, color=red](J,K)
%		\tkzDefMidPoint(B,C)\tkzGetPoint{D}
%		\tkzDefMidPoint(J,K)\tkzGetPoint{L}
%		\draw[<->, >=latex] (A) .. controls +(right:1cm) and +(left:2cm) .. (I) node[sloped,above,pos=0.5]{\scriptsize Angles homologues} ;
%		\draw[<->, >=latex] (D) .. controls +(right:8cm) and +(down:4cm) .. (L) node[ below right,pos=0.6]{\scriptsize Côtés homologues} ;
%	\end{tikzpicture}
%	\caption{$ABC$ et $IJK$ sont égaux.} \label{chapitre05:trianglesegaux}
%\end{figure}  
%\begin{align*}
%	& {\color{red} \widehat{A}= \widehat{I} }
%	&&  {\color{ForestGreen} \widehat{B}= \widehat{J} } 
%	&& {\color{Blue} \widehat{C}= \widehat{K} } \\
%	& {\color{red} BC=JK }
%	&&  {\color{ForestGreen}  AC=IK } 
%	&& {\color{Blue} AB=IJ} 
%\end{align*}
%
%Il n'est pas nécessaire de vérifier les six égalités pour prouver que les triangles sont égaux et que les 6 égalités sont vraies.  
%Il suffit de vérifier que l'on est dans l'une des trois situations suivantes.

\begin{figure}[hb]
	\begin{tikzpicture} 
		\tkzDefPoint(0,0){A}
		%\tkzLabelPoint[above left](A){$A$}
		\tkzDefShiftPoint[A](-135:2.5){B}  
		\tkzDefTriangle[two angles = 80 and 60](A,B)
		\tkzGetPoint{C}
		%\tkzLabelPoint[left](B){$B$}
		%\tkzLabelPoint[right](C){$C$}
		\tkzDrawPolygon(A,B,C)
		\tkzDrawPoints(A,B,C)
		\tkzDefShiftPoint[A](6.0,-2.0){J}  
		\tkzDefShiftPoint[J](110:2.5){I}
		\tkzDefTriangle[two angles = -80  and -60](I,J)
		\tkzGetPoint{K}
		\tkzDrawPolygon(I,J,K)
		\tkzDrawPoints(I,J,K) 
		%\tkzLabelPoint[right](I){$I$}
		%\tkzLabelPoint[below left](J){$J$}
		%\tkzLabelPoint[below](K){$K$}
		\tkzMarkSegments[mark=s|, color=black, thick](B,C K,J)
		\tkzMarkSegments[mark=s||, color=black, thick](B,A I,J)
		\tkzMarkSegments[mark=x,color=black, thick](A,C K,I)
		%\tkzMarkAngle[mark=|,mkcolor=black, thick, size=0.6, fill=red!20, opacity=0.6](B,A,C)
		%\tkzMarkAngle[mark=||,mkcolor=black, thick,size=0.6, fill=blue!20, opacity=0.6](J,K,I)
		%\tkzMarkAngle[mark=||,mkcolor=black, thick,size=0.6, fill=blue!20, opacity=0.6](A,C,B)
		%\tkzMarkAngle[mark=|,mkcolor=black, thick,size=0.6, fill=red!20, opacity=0.6](K,I,J)
		%\tkzMarkAngle[mark=s,mkcolor=black, thick,size=0.6, fill=green!20, opacity=0.6](C,B,A)
		%\tkzMarkAngle[mark=s,mkcolor=black, thick,size=0.6, fill=green!20, opacity=0.6](I,J,K) 
		\tkzDrawSegment[line width=1pt ](B,C)
		\tkzDrawSegment[line width=1pt ](J,K)
		\tkzDefMidPoint(B,C)\tkzGetPoint{D}
		\tkzDefMidPoint(J,K)\tkzGetPoint{L}
	\end{tikzpicture}
\caption{Critère CCC : Si deux triangles ont leurs trois côtés respectivement égaux, alors ils sont égaux.}
\end{figure}    
	\begin{figure}[hb]
		\begin{tikzpicture}
			\tkzDefPoint(0,0){A}
			%\tkzLabelPoint[above left](A){$A$}
			\tkzDefShiftPoint[A](-125:2.5){B}  
			\tkzDefTriangle[two angles = 80 and 60](A,B)
			\tkzGetPoint{C}
			%\tkzLabelPoint[left](B){$B$}
			%\tkzLabelPoint[right](C){$C$}
			\tkzDrawPolygon(A,B,C)
			\tkzDrawPoints(A,B,C)
			\tkzDefShiftPoint[A](6.0,-2.0){J}  
			\tkzDefShiftPoint[J](110:2.5){I}
			\tkzDefTriangle[two angles = -80  and -60](I,J)
			\tkzGetPoint{K}
			\tkzDrawPolygon(I,J,K)
			\tkzDrawPoints(I,J,K) 
			%\tkzLabelPoint[right](I){$I$}
			%\tkzLabelPoint[below left](J){$J$}
			%\tkzLabelPoint[below](K){$K$}
			%\tkzMarkSegments[mark=x,mkcolor=black, thick](B,C K,J)
			\tkzMarkSegments[mark=s|,color=black, thick](B,A I,J)
			\tkzMarkSegments[mark=s||,color=black, thick](A,C K,I)
			\tkzMarkAngle[mark=s,color=black, thick, size=0.6, fill=red!20, opacity=0.6](B,A,C)
			%\tkzMarkAngle[mark=||,mkcolor=black, thick,size=0.6, fill=blue!20, opacity=0.6](J,K,I)
			%\tkzMarkAngle[mark=||,mkcolor=black, thick,size=0.6, fill=blue!20, opacity=0.6](A,C,B)
			\tkzMarkAngle[mark=s,mkcolor=black, thick,size=0.6, fill=red!20, opacity=0.6](K,I,J)
			%\tkzMarkAngle[mark=s,mkcolor=black, thick,size=0.6, fill=green!20, opacity=0.6](C,B,A)
			%\tkzMarkAngle[mark=s,mkcolor=black, thick,size=0.6, fill=green!20, opacity=0.6](I,J,K) 
			\tkzDrawSegment[line width=1pt ](B,C)
			\tkzDrawSegment[line width=1pt ](J,K)
			\tkzDefMidPoint(B,C)\tkzGetPoint{D}
			\tkzDefMidPoint(J,K)\tkzGetPoint{L}
		\end{tikzpicture}
	\caption{Critère CAC : Si deux triangles ont un angle égal compris entre deux côtés respectivement égaux, alors ils sont égaux.}
	\end{figure} 
%\end{postulat} 

%\begin{postulat}[Critère ACA] Si deux triangles ont un côté égal adjacent à deux angles respectivement égaux, alors ils sont égaux.
	\begin{figure}[hb]
		\begin{tikzpicture}
			\tkzDefPoint(0,0){A}
			%\tkzLabelPoint[above left](A){$A$}
			\tkzDefShiftPoint[A](-135:2.5){B}  
			\tkzDefTriangle[two angles = 80 and 60](A,B)
			\tkzGetPoint{C}
			%\tkzLabelPoint[left](B){$B$}
			%\tkzLabelPoint[right](C){$C$}
			\tkzDrawPolygon(A,B,C)
			\tkzDrawPoints(A,B,C)
			\tkzDefShiftPoint[A](6.0,-2.0){J}  
			\tkzDefShiftPoint[J](110:2.5){I}
			\tkzDefTriangle[two angles = -80  and -60](I,J)
			\tkzGetPoint{K}
			\tkzDrawPolygon(I,J,K)
			\tkzDrawPoints(I,J,K) 
			%\tkzLabelPoint[right](I){$I$}
			%\tkzLabelPoint[below left](J){$J$}
			%\tkzLabelPoint[below](K){$K$}
			\tkzMarkSegments[mark=s||,color=black, thick](B,C K,J)
			%\tkzMarkSegments[mark=x,mkcolor=black, thick](B,A I,J)
			%\tkzMarkSegments[mark=s|,mkcolor=black, thick](A,C K,I)
			%\tkzMarkAngle[mark=||,mkcolor=black, thick, size=0.6, fill=red!20, opacity=0.6](B,A,C)
			\tkzMarkAngle[mark=|,mkcolor=black, thick,size=0.6, fill=blue!20, opacity=0.6](J,K,I)
			\tkzMarkAngle[mark=|,mkcolor=black, thick,size=0.6, fill=blue!20, opacity=0.6](A,C,B)
			%\tkzMarkAngle[mark=||,mkcolor=black, thick,size=0.6, fill=red!20, opacity=0.6](K,I,J)
			\tkzMarkAngle[mark=s,mkcolor=black, thick,size=0.6, fill=green!20, opacity=0.6](C,B,A)
			\tkzMarkAngle[mark=s,mkcolor=black, thick,size=0.6, fill=green!20, opacity=0.6](I,J,K) 
			\tkzDrawSegment[line width=1pt ](B,C)
			\tkzDrawSegment[line width=1pt ](J,K)
			\tkzDefMidPoint(B,C)\tkzGetPoint{D}
			\tkzDefMidPoint(J,K)\tkzGetPoint{L}
		\end{tikzpicture}
	\caption{Critère ACA Si deux triangles ont un côté égal adjacent à deux angles respectivement égaux, alors ils sont égaux.}
	\end{figure}
%\end{postulat}  
\begin{remark}
	Cas RHC (Rectangle-Hypoténuse-Côté). Deux triangles rectangles, qui ont même longueur d'hypoténuse et une même longueur d'un côté de l'angle droit sont égaux.
\end{remark}
%\begin{remark} Il n'y a pas de critère ACC
%	% Tracer deux triangles non égaux, ayant un angle de \ang{40}, adjacent à un côté de \numprint{4.5}~\si{\centi\meter} et opposé à un côté de \numprint{3.5}~\si{\centi\meter} de longueur.
%	\begin{figure}[h]\centering
%		\begin{tikzpicture}[scale=1.0]
%			\tkzDefPoint(0,0){A}
%			\tkzDefShiftPoint[A](40:5){B}
%			\tkzDefMidPoint(A,B)\tkzGetPoint{I}
%			\tkzDefShiftPoint[I](-0.2,0.2){I'}
%			\node[inner sep=0pt, opacity=1,   rotate=-15] (carte) at (I') {
%				\includegraphics[ width=0.25\linewidth]{fishing-pole2}
%			}; 
%			\tkzDefShiftPoint[A](0:6){C}
%			%\tkzDrawPoint(B)
%			\tkzDrawLine[add= 0.2 and 0.25, dashed](A,C)
%			\tkzDrawLine[add= 0 and 0.15, dashed](A,B)
%			\tkzLabelAngle[pos=1.5](C,A,B){\large\ang{40}} 
%			\tkzMarkAngle[size=1.0, opacity=0.5,mark=none](C,A,B)
%		\end{tikzpicture}
%		\caption{}
%	\end{figure} 
%\end{remark}  




\begin{definition} Deux triangles sont \textbf{semblables}   lorsqu'ils ont leurs angles égaux deux à deux et leurs côtés \textbf{proportionnels}. 
\end{definition}

\begin{figure}[h]\centering
	\begin{tikzpicture}[scale=0.8]
		\tkzDefPoint(0,0){A}
		\tkzLabelPoint[above left](A){$A$}
		\tkzDefShiftPoint[A](-135:2.5){B}  
		\tkzDefTriangle[two angles = 80 and 60](A,B)
		\tkzGetPoint{C}
		\tkzLabelPoint[left](B){$B$}
		\tkzLabelPoint[right](C){$C$}
		\tkzDrawPolygon(A,B,C)
		\tkzDrawPoints(A,B,C)
		\tkzDefShiftPoint[A](4,-1.0){J}  
		\tkzDefShiftPoint[J](10:6){I}
		\tkzDefTriangle[two angles = -60 and -40](I,J)
		\tkzGetPoint{K}
		\tkzDrawPolygon(I,J,K)
		\tkzDrawPoints(I,J,K) 
		\tkzLabelPoint[right](I){$I$}
		\tkzLabelPoint[left](J){$J$}
		\tkzLabelPoint[right](K){$K$}
		\tkzMarkAngle[mark=|,mkcolor=black, thick, size=0.6, fill=red!20, opacity=0.6](B,A,C)
		\tkzMarkAngle[mark=|,mkcolor=black, thick,size=0.6, fill=red!20, opacity=0.6](J,K,I)
		\tkzMarkAngle[mark=||,mkcolor=black, thick,size=0.6, fill=blue!20, opacity=0.6](A,C,B)
		\tkzMarkAngle[mark=||,mkcolor=black, thick,size=0.6, fill=blue!20, opacity=0.6](I,J,K) 
		\tkzMarkAngle[mark=s,mkcolor=black, thick,size=0.6, fill=green!20, opacity=0.6](C,B,A)
		\tkzMarkAngle[mark=s,mkcolor=black, thick,size=0.6, fill=green!20, opacity=0.6](K,I,J)
		\tkzDrawSegment[line width=2pt, color=red](B,C)
		\tkzDrawSegment[line width=2pt, color=red](I,J)
		\tkzDefMidPoint(B,C)\tkzGetPoint{D}
		\tkzDefMidPoint(I,J)\tkzGetPoint{L}
		\draw[<->, >=latex] (A) .. controls +(right:1cm) and +(left:2cm) .. (K) node[sloped,above,pos=0.5]{\scriptsize Angles homologues} ;
		\draw[<->, >=latex] (D) .. controls +(right:8cm) and +(down:3cm) .. (L) node[ below right,pos=0.6]{\scriptsize Côtés homologues} ;
	\end{tikzpicture}
	\caption{les trianges $ABC$ et $IJK$ sont semblables.}
\end{figure} 
Les \textbf{angles correspondants} sont égaux :
\[{\color{red} \widehat{A}= \widehat{K} }\qquad {\color{blue} \widehat{B}= \widehat{J} }\qquad {\color{ForestGreen} \widehat{C}= \widehat{I} }\]
Les \textbf{côtés correspondants} sont proportionnels  :    
%\begin{table}[h]
%	\renewcommand{\arraystretch}{2}%
%	\begin{tabularx}{0.9\linewidth}{|X?c|c|c|} \hline 
%		\tikz[remember picture,baseline =(C.base)]\node (C) at (-1cm,0.5cm) {};	{\color{HotPink} Côtés  du triangle $ABC$} &  \emph{\color{blue}$AB$}&  \emph{\color{red}$BC$} & \emph{$AC$} \tikz[remember picture]   \node  (A) at (-1cm,0.5cm) {}; \\   \hline
%		\tikz[remember picture,baseline =(D.base)]\node (D) at (-1cm,0.5cm) {};	{\color{ForestGreen} Côtés  du triangle $IJK$} &  \emph{\color{blue}$JK$}&  \emph{\color{red}$IJ$}& \emph{$IK$} \tikz[remember picture,baseline =(B.base)]\node (B) at (-1cm,0.5cm) {}; \\ \hline
%	\end{tabularx}
%	\tikz[overlay,remember picture]
%	\draw[ ->,>=latex,line width=1.8pt ](A) .. controls +(+1.5cm,.1cm) and +(+1.5cm,-.0cm)..
%	node[rectangle,fill=white,draw,circle ]{\Large$\times k$} (B);
%	%	\tikz[overlay,remember picture]
%	%	\draw[ ->,>=latex,line width=1.8pt ](C) .. controls +(-1.5cm,.1cm) and +(-1.5cm,-.0cm)..
%	%	node[rectangle,fill=white,draw,circle ]{\Large$\times k$} (D);
%	\caption{ Si $k> 1$, le triangle $IJK$ est un \emph{agrandissement} de $ABC$.
%		Si $k< 1$, le triangle $IJK$ est une \emph{réduction} de $ABC$.}
%\end{table}
\[ 	{\color{blue} \frac{JK}{AB} } = {\color{red} \frac{IJ}{BC} } = {\color{ForestGreen} \frac{IK}{AC} } \color{black} =k \]
%\begin{remark}
%	Pour un rapport d'agrandissement (ou réduction) $k$.
%	\begin{itemize}
%		\item Les longueurs sont multipliées par $k$.
%		\item Les aires sont multipliées par $k^2$.
%	\end{itemize} 
%\end{remark}  
\begin{postulat}[Critère de similitude CCC] Si les longueurs des 3 côtés d'un triangle $T_1$ sont \textbf{proportionnelles} aux longueurs respectives des 3 côtés d'un triangle $T_2$, alors les deux triangles sont semblables.
\end{postulat}

\begin{postulat}[Critère CAC-Semblable] Si deux triangles $T_1$ et $T_2$ ont un angle égal compris entre 2 côtés respectivement proportionnels, alors les deux triangles sont semblables.  
\end{postulat} 
 
\begin{postulat}[Critère de similitude AA] Si 2 angles d'un triangle $T_1$ sont \textbf{respectivement égaux} à 2 angles d'un triangle $T_2$. Alors les deux triangles sont semblables.  
\end{postulat} 

%----------------------------------------------------------------------------------------
%	EXERCICES
%----------------------------------------------------------------------------------------

\newpage 
\loadgeometry{widelayout}  
\pagestyle{scrheadings} % Print headers again  
\KOMAoptions{
	headwidth=textwithmarginpar,
	headsepline=true,
	headsepline= 0.5pt,
	footwidth=textwithmarginpar,
} 
\onehalfspacing
\subsection{Exercices : calculs algébriques et géométrie}
 
\begin{exercise}[Triangles égaux : à l'oral]
	
	Si possible, démontrer pour chaque cas que les triangles sont égaux. Les figures ne sont pas à l'échelle. Utiliser la même couleur pour indiquer les angles et côtés homologues.
	
	\noindent\parbox[position=t]{0.33\textwidth}{
		%\begin{figure}[H]\centering
		\begin{tikzpicture}[scale=0.75]
			\tkzText(0.5,2.8){$(A)$}
			\tkzSetUpPoint[size=3, fill=black]
			\tkzfctset{tan style/.style={->,>=latex, color=red, line width=1.2pt}}   
			\tkzSetUpLine[line width= 1.2pt, add=.5 and .5] 
			\tkzDefPoint(1,0){A} 
			\tkzDefShiftPoint[A](-5:3.3){B}
			\tkzDefTriangle[two angles=90 and 35](A,B)\tkzGetPoint{C}  
			\tkzDrawPolygon[line width=1.2pt](A,B,C)
			%		\tkzLabelPoints[font=\scriptsize](A,B,C)
			\tkzMarkRightAngle[](B,A,C)
			%\tkzMarkAngle[mark=none,-, >=latex, size =0.5](C,B,A)
			%\tkzMarkAngle[mark=none,-, >=latex, size =0.5](A,C,B)
			%\tkzLabelAngle[font=\scriptsize,pos=.8](C,B,A){$\ang{70}$}
			%\tkzLabelAngle[font=\scriptsize,pos=.8](A,C,B){$\ang{55}$}
			%\tkzLabelSegment[above, font=\scriptsize](A,C){$10$} 
			\tkzLabelSegment[below, font=\scriptsize](A,B){\SI{7}{\centi\meter}}
			\tkzLabelSegment[above,sloped,  font=\scriptsize](C,B){\SI{9}{\centi\meter}}
			\begin{scope}[xshift=5cm] 
				\tkzDefPoint(0,0){A}
				\tkzDefShiftPoint[A](+15:2.5){B} 
				\tkzDefTriangle[two angles=90 and 35](A,B)\tkzGetPoint{C}  
				
				\tkzDrawPolygon[line width=1.2pt](A,B,C)
				%	\tkzLabelPoints[font=\scriptsize](A,B,C)
				\tkzMarkRightAngle[](B,A,C)
				%\tkzMarkAngle[mark=none,-, >=latex, size =0.5](C,B,A)
				%\tkzMarkAngle[mark=none,-, >=latex, size =0.5](A,C,B)
				%\tkzLabelAngle[font=\scriptsize,pos=.8](C,B,A){$\ang{70}$}
				%\tkzLabelAngle[font=\scriptsize,pos=.8](A,C,B){$\ang{55}$}
				%\tkzLabelSegment[above, font=\scriptsize](A,C){$10$} 
				\tkzLabelSegment[left, font=\scriptsize](A,C){\SI{7}{\centi\meter}}
				\tkzLabelSegment[above right, font=\scriptsize](C,B){\SI{9}{\centi\meter}}
			\end{scope}
		\end{tikzpicture}   
	}
%	\hfill 
	\parbox[position=t]{0.33\textwidth}{
		%\begin{figure}[H]\centering
		\begin{tikzpicture}[scale=0.75]
			\tkzText(0.5,2){$(B)$}
			\tkzSetUpPoint[size=3, fill=black]
			\tkzfctset{tan style/.style={->,>=latex, color=red, line width=1.2pt}}   
			\tkzSetUpLine[line width= 1.2pt, add=.5 and .5] 
			\tkzDefPoint(0,0){A}
			\tkzDefPoint(3,0){B}
			\tkzDefShiftPoint[A](-26:3.1){B}
			\tkzDefShiftPoint[A](32:4.2){C} 
			\tkzDrawPolygon[line width=1.2pt](A,B,C)
			%		\tkzLabelPoints[font=\scriptsize](A,B,C)
			\tkzMarkAngle[mark=none,-, >=latex, size =0.5](C,B,A)
			\tkzMarkAngle[mark=none,-, >=latex, size =0.5](A,C,B)
			\tkzLabelAngle[font=\scriptsize,pos=.9](C,B,A){$\ang{70}$}
			\tkzLabelAngle[font=\scriptsize,pos=1.0](A,C,B){$\ang{60}$}
			\tkzLabelSegment[above left, font=\scriptsize](A,C){\SI{10}{\centi\meter}}
			%		\tkzDefPointWith[orthogonal,K= -0.65](C,A)
			%		\tkzGetPoint{D}
			%		\tkzDefLine[parallel=through D](A,B) \tkzGetPoint{E}
			%		\tkzDrawLine(D,E)
			%		\tkzDrawLine(A,B)
			%		\tkzLabelLine[pos=1.5,  right](A,B){$(\Delta_1)$}
			%		\tkzLabelLine[pos=1.5,  right](D,E){$(\Delta_2)$}
			%		\tkzMarkRightAngle(D,C,A)
			%		
			%		\tkzDrawSegments[line width= 1.2pt](A,C C,D)
			%		\tkzMarkAngle[mark=none,-, >=latex, size =0.7](B,A,C)
			%		\tkzLabelAngles[pos=1.3,  ](B,A,C){$a$}
			%		\tkzMarkAngle[mark=none,-, >=latex, size =0.7](C,D,E)
			%		\tkzLabelAngles[pos=1.3, ](C,D,E){$b$} 
			\begin{scope}[xshift=4cm] 
				\tkzDefPoint(0,0){A}
				\tkzDefPoint(3,0){B}
				\tkzDefShiftPoint[A](-5:4.1){B}
				\tkzDefShiftPoint[A](48:3.3){C} 
				\tkzDrawPolygon[line width=1.2pt](A,B,C)
				%	\tkzLabelPoints[font=\scriptsize](A,B,C)
				\tkzMarkAngle[mark=none,-, >=latex, size =0.5](B,A,C)
				\tkzMarkAngle[mark=none,-, >=latex, size =0.5](A,C,B)
				\tkzLabelAngle[font=\scriptsize,pos=.9](B,A,C){$\ang{50}$}
				\tkzLabelAngle[font=\scriptsize,pos=.8](A,C,B){$\ang{70}$} 
				\tkzLabelSegment[below, font=\scriptsize](A,B){\SI{10}{\centi\meter}}
			\end{scope}
		\end{tikzpicture}   
	}
%\hfill
\parbox[position=t]{0.33\textwidth}{
		%\begin{figure}[H]\centering
		\begin{tikzpicture}[scale=0.75]
			\tkzText(0.5,2.1){$(C)$}
			\tkzSetUpPoint[size=3, fill=black]
			\tkzfctset{tan style/.style={->,>=latex, color=red, line width=1.2pt}}   
			\tkzSetUpLine[line width= 1.2pt, add=.5 and .5] 
			\tkzDefPoint(0.5,0){A} 
			\tkzDefShiftPoint[A](0:3.7){B}
			\tkzDefTriangle[two angles=-25 and -110](A,B)\tkzGetPoint{C}  
			\tkzDrawPolygon[line width=1.2pt](A,B,C)
			%		\tkzLabelPoints[font=\scriptsize](A,B,C)
			%		\tkzMarkRightAngle[](B,A,C)
			\tkzMarkAngle[mark=none,-, >=latex, size =0.8](C,A,B)
			\tkzMarkAngle[mark=none,-, >=latex, size =0.4](A,B,C)
			\tkzMarkAngle[mark=none,-, >=latex, size =0.6](B,C,A)
			\tkzLabelAngle[font=\scriptsize,pos=1.6](C,A,B){$\ang{20}$}
			\tkzLabelAngle[font=\scriptsize,pos=0.7](A,B,C){$\ang{130}$}
			\tkzLabelAngle[font=\scriptsize,pos=1.2](B,C,A){$\ang{30}$}
			%\tkzLabelSegment[above, font=\scriptsize](A,C){ \SI{13}{\centi\meter}} 
			%		\tkzLabelSegment[above, font=\scriptsize](A,B){\SI{8}{\centi\meter}}
			%		\tkzLabelSegment[above right, font=\scriptsize, sloped](C,B){\SI{9}{\centi\meter}}
			\begin{scope}[xshift=4.5cm, yshift= -0.5cm] 
				\tkzDefPoint(0.5,0){A} 
				\tkzDefShiftPoint[A](-5:3.7){B}
				\tkzDefTriangle[two angles=120 and 30](A,B)\tkzGetPoint{C}  
				\tkzDrawPolygon[line width=1.2pt](A,B,C)
				%		\tkzLabelPoints[font=\scriptsize](A,B,C)
				%		\tkzMarkRightAngle[](B,A,C)
				\tkzMarkAngle[mark=none,-, >=latex, size =0.5](B,A,C)
				\tkzMarkAngle[mark=none,-, >=latex, size =0.8](C,B,A)
				\tkzMarkAngle[mark=none,-, >=latex, size =1.0](A,C,B)
				\tkzLabelAngle[font=\scriptsize,pos=0.8](B,A,C){$\ang{130}$}
				\tkzLabelAngle[font=\scriptsize,pos=1.7](C,B,A){$\ang{20}$}
				\tkzLabelAngle[font=\scriptsize,pos=1.5](A,C,B){$\ang{30}$}
				%	\tkzLabelSegment[left, font=\scriptsize](A,C){\SI{9}{\centi\meter}} 
				%			\tkzLabelSegment[below  , font=\scriptsize](A,B){\SI{8}{\centi\meter}}
				%		\tkzLabelSegment[right, font=\scriptsize](C,B){\SI{8}{\centi\meter}}
			\end{scope}
		\end{tikzpicture}   
	}


\noindent\parbox[position=t]{0.4\textwidth}{
		%\begin{figure}[H]\centering
		\begin{tikzpicture}[scale=0.75]
			\tkzText(0.5,2.8){$(D)$}
			\tkzSetUpPoint[size=3, fill=black]
			\tkzfctset{tan style/.style={->,>=latex, color=red, line width=1.2pt}}   
			\tkzSetUpLine[line width= 1.2pt, add=.5 and .5] 
			\tkzDefPoint(0.5,0){A} 
			\tkzDefShiftPoint[A](0:4.1){B}
			\tkzDefTriangle[two angles=25 and 125](A,B)\tkzGetPoint{C}  
			\tkzDrawPolygon[line width=1.2pt](A,B,C)
			%		\tkzLabelPoints[font=\scriptsize](A,B,C)
			%		\tkzMarkRightAngle[](B,A,C)
			%		\tkzMarkAngle[mark=none,-, >=latex, size =0.6](B,A,C)
			\tkzMarkAngle[mark=none,-, >=latex, size =0.4](C,B,A)
			%		\tkzMarkAngle[mark=none,-, >=latex, size =0.5](A,C,B)
			%		\tkzLabelAngle[font=\scriptsize,pos=1.2](B,A,C){$\ang{20}$}
			\tkzLabelAngle[font=\scriptsize,pos=0.7](C,B,A){$\ang{135}$}
			%		\tkzLabelAngle[font=\scriptsize,pos=.8](A,C,B){$\ang{45}$}
			%		\tkzLabelSegment[above, font=\scriptsize](A,C){ \SI{13}{\centi\meter}} 
			\tkzLabelSegment[below, font=\scriptsize](A,B){\SI{15}{\centi\meter}}
			\tkzLabelSegment[above,sloped, font=\scriptsize](C,B){\SI{8}{\centi\meter}}
			\begin{scope}[xshift=4.5cm, yshift=-1cm] 
				\tkzDefPoint(0.5,0){A} 
				\tkzDefShiftPoint[A](+20:3.1){B}
				\tkzDefTriangle[two angles=35 and 115](A,B)\tkzGetPoint{C}  
				\tkzDrawPolygon[line width=1.2pt](A,B,C)
				%		\tkzLabelPoints[font=\scriptsize](A,B,C)
				%		\tkzMarkRightAngle[](B,A,C)
				\tkzMarkAngle[mark=none,-, >=latex, size =0.7](B,A,C)
				%			\tkzMarkAngle[mark=none,-, >=latex, size =0.5](C,B,A)
				\tkzMarkAngle[mark=none,-, >=latex, size =0.7](A,C,B)
				\tkzLabelAngle[font=\scriptsize,pos=1.5](B,A,C){$\ang{20}$}
				%\tkzLabelAngle[font=\scriptsize,pos=.8](C,B,A){$\ang{35}$}
				\tkzLabelAngle[font=\scriptsize,pos=1.1](A,C,B){$\ang{25}$}
				%		\tkzLabelSegment[left, font=\scriptsize](A,C){$9$} 
				\tkzLabelSegment[below right, font=\scriptsize](A,B){\SI{15}{\centi\meter}}
				\tkzLabelSegment[right, font=\scriptsize](C,B){\SI{8}{\centi\meter}}
			\end{scope}
		\end{tikzpicture}   
	}
%	\hfill 
 \parbox[position=t]{0.30\textwidth}{
		%\begin{figure}[H]\centering
		\begin{tikzpicture}[scale=0.75]
			\tkzText(0.5,2.5){$(E)$}
			\tkzSetUpPoint[size=3, fill=black]
			\tkzfctset{tan style/.style={->,>=latex, color=red, line width=1.2pt}}   
			\tkzSetUpLine[line width= 1.2pt, add=.5 and .5] 
			\tkzDefPoint(0.5,0){A} 
			\tkzDefShiftPoint[A](-20:2.2){B}
			\tkzDefTriangle[two angles=90 and 55](A,B)\tkzGetPoint{C}  
			\tkzDrawPolygon[line width=1.2pt](A,B,C)
			%		\tkzLabelPoints[font=\scriptsize](A,B,C)
			\tkzMarkRightAngle[](B,A,C)
			%		\tkzMarkAngle[mark=none,-, >=latex, size =0.6](B,A,C)
			%		\tkzMarkAngle[mark=none,-, >=latex, size =0.4](C,B,A)
			%		\tkzMarkAngle[mark=none,-, >=latex, size =0.5](A,C,B)
			%		\tkzLabelAngle[font=\scriptsize,pos=1.2](B,A,C){$\ang{20}$}
			%		\tkzLabelAngle[font=\scriptsize,pos=0.6](C,B,A){$\ang{70}$}
			%		\tkzLabelAngle[font=\scriptsize,pos=.8](A,C,B){$\ang{55}$}
			\tkzLabelSegment[left, font=\scriptsize](A,C){ \SI{9}{\centi\meter}} 
			%	\tkzLabelSegment[below, font=\scriptsize](A,B){\SI{10}{\centi\meter}}
			\tkzLabelSegment[above right, font=\scriptsize](C,B){\SI{12}{\centi\meter}}
			\begin{scope}[xshift=3.5cm, yshift=-1cm] 
				\tkzDefPoint(0.5,0){A} 
				\tkzDefShiftPoint[A](30:3.1){B}
				\tkzDefTriangle[two angles=35 and 90](A,B)\tkzGetPoint{C}  
				\tkzDrawPolygon[line width=1.2pt](A,B,C)
				%		\tkzLabelPoints[font=\scriptsize](A,B,C)
				%					\tkzMarkRightAngle[](B,A,C)
				\tkzMarkRightAngle[](C,B,A) 
				%			\tkzMarkAngle[mark=none,-, >=latex, size =0.6](B,A,C)
				%			\tkzMarkAngle[mark=none,-, >=latex, size =0.5](C,B,A)
				%			\tkzMarkAngle[mark=none,-, >=latex, size =0.6](A,C,B)
				%			\tkzLabelAngle[font=\scriptsize,pos=0.9](B,A,C){$\ang{55}$}
				%\tkzLabelAngle[font=\scriptsize,pos=.8](C,B,A){$\ang{35}$}
				%			\tkzLabelAngle[font=\scriptsize,pos=0.9](A,C,B){$\ang{70}$}
				\tkzLabelSegment[above left, font=\scriptsize](A,C){\SI{12}{\centi\meter}} 
				% \tkzLabelSegment[below  , font=\scriptsize](A,B){\SI{10}{\centi\meter}}
				\tkzLabelSegment[above right, font=\scriptsize](C,B){\SI{8}{\centi\meter}}
			\end{scope}
		\end{tikzpicture}   
	}
%	\hfill 
	\parbox[position=t]{0.30\textwidth}{
		%\begin{figure}[H]\centering
		\begin{tikzpicture}[scale=0.75]
			\tkzText(0.5,2.5){$(F)$}
			\tkzSetUpPoint[size=3, fill=black]
			\tkzfctset{tan style/.style={->,>=latex, color=red, line width=1.2pt}}   
			\tkzSetUpLine[line width= 1.2pt, add=.5 and .5] 
			\tkzDefPoint(0.5,0){A} 
			\tkzDefShiftPoint[A](-20:2.5){B}
			\tkzDefTriangle[two angles=80 and 50](A,B)\tkzGetPoint{C}  
			\tkzDrawPolygon[line width=1.2pt](A,B,C)
			%		\tkzLabelPoints[font=\scriptsize](A,B,C)
			%		\tkzMarkRightAngle[](B,A,C)
			%		\tkzMarkAngle[mark=none,-, >=latex, size =0.6](B,A,C)
			\tkzMarkAngle[mark=none,-, >=latex, size =0.4](C,B,A)
			\tkzMarkAngle[mark=none,-, >=latex, size =0.5](A,C,B)
			%		\tkzLabelAngle[font=\scriptsize,pos=1.2](B,A,C){$\ang{20}$}
			\tkzLabelAngle[font=\scriptsize,pos=1.0](C,B,A){$\ang{50}$}
			\tkzLabelAngle[font=\scriptsize,pos=1.0](A,C,B){$\ang{40}$}
			\tkzLabelSegment[left, font=\scriptsize](A,C){ \SI{9}{\centi\meter}} 
			%	\tkzLabelSegment[below, font=\scriptsize](A,B){\SI{10}{\centi\meter}}
			%		\tkzLabelSegment[above right, font=\scriptsize](C,B){\SI{12}{\centi\meter}}
			\begin{scope}[xshift=3.5cm, yshift=+2cm] 
				\tkzDefPoint(0.5,0){A} 
				\tkzDefShiftPoint[A](-60:3.1){B}
				\tkzDefTriangle[two angles=90 and 35](A,B)\tkzGetPoint{C}  
				\tkzDrawPolygon[line width=1.2pt](A,B,C)
				%		\tkzLabelPoints[font=\scriptsize](A,B,C)
				\tkzMarkRightAngle[](B,A,C)
				%			\tkzMarkRightAngle[](C,B,A) 
				%			\tkzMarkAngle[mark=none,-, >=latex, size =0.6](B,A,C)
				%			\tkzMarkAngle[mark=none,-, >=latex, size =0.5](C,B,A)
				%			\tkzMarkAngle[mark=none,-, >=latex, size =0.6](A,C,B)
				%			\tkzLabelAngle[font=\scriptsize,pos=0.9](B,A,C){$\ang{55}$}
				%\tkzLabelAngle[font=\scriptsize,pos=.8](C,B,A){$\ang{35}$}
				%			\tkzLabelAngle[font=\scriptsize,pos=0.9](A,C,B){$\ang{70}$}
				\tkzLabelSegment[above left, font=\scriptsize](A,C){\SI{9}{\centi\meter}} 
				% \tkzLabelSegment[below  , font=\scriptsize](A,B){\SI{10}{\centi\meter}}
				\tkzLabelSegment[above right, font=\scriptsize](C,B){\SI{12}{\centi\meter}}
			\end{scope}
		\end{tikzpicture}   
	} 
\end{exercise}
\vfill  
\begin{exercise}[Révision triangles semblables] \hfill 
	
	
\noindent\parbox{0.4\linewidth}{\centering
\begin{tikzpicture}[scale=1.0] 
	\tkzSetUpPoint[size=3, fill=black]
	\tkzfctset{tan style/.style={->,>=latex, color=red, line width=1.2pt}}   
	\tkzSetUpLine[line width= 1.2pt, add=.5 and .5] 
	\tkzDefPoint(0.0,0){A} 
	\tkzDefShiftPoint[A](-25:3.5){B}
	\tkzDefTriangle[two angles=35 and 70](A,B)\tkzGetPoint{C}  
	\tkzDrawPolygon[line width=1.2pt](A,B,C)
	\tkzLabelPoints[above,font=\scriptsize](C)
	\tkzLabelPoints[right,font=\scriptsize](B)
	\tkzLabelPoints[left,font=\scriptsize](A)
	%	\tkzMarkRightAngle[](B,A,C)
	%			\tkzMarkAngle[mark=||,-, >=latex, size =0.5](B,A,C)
	\tkzMarkAngle[mark=|,-, >=latex, size =0.7](C,B,A)
	\tkzMarkAngle[mark=||,-, >=latex, size =0.5](A,C,B)
	%		\tkzLabelAngle[font=\scriptsize,pos=.8](B,A,C){$\ang{80}$}
	%		\tkzLabelAngle[font=\scriptsize,pos=.8](C,B,A){$\ang{30}$}
	%		\tkzLabelAngle[font=\scriptsize,pos=.8](A,C,B){$\ang{70}$} 
	\tkzLabelSegment[above, sloped , font=\scriptsize](A,C){\SI{14}{\meter}} 
	\tkzLabelSegment[below, sloped , font=\scriptsize](A,B){$x$}
	\tkzLabelSegment[below,sloped ,  font=\scriptsize](C,B){\SI{12}{\meter}}
	\begin{scope}[xshift=4.5cm,yshift=-2cm ] 
		\tkzDefPoint(0.0,0){D} 
		\tkzDefShiftPoint[D](60:2.5){E}
		\tkzDefTriangle[two angles=30 and 75](D,E)\tkzGetPoint{F}  
		\tkzDrawPolygon[line width=1.2pt](D,E,F)
		\tkzLabelPoints[below left,font=\scriptsize](D)
		\tkzLabelPoints[below right,font=\scriptsize](E)
		\tkzLabelPoints[above,font=\scriptsize](F) 
		%			\tkzMarkAngle[mark=none,-, >=latex, size =0.5](E,D,F)
		\tkzMarkAngle[mark=|,-, >=latex, size =0.5](F,E,D)
		\tkzMarkAngle[mark=||,-, >=latex, size =0.5](D,F,E)
		%			\tkzLabelAngle[font=\scriptsize,pos=.8](E,D,F){$\ang{70}$}
		%			\tkzLabelAngle[font=\scriptsize,pos=.8](F,E,D){$\ang{30}$}
		%			\tkzLabelAngle[font=\scriptsize,pos=.8](D,F,E){$\ang{80}$}
		\tkzLabelSegment[below,sloped , font=\scriptsize](D,E){\SI{10}{\centi\meter}}
		\tkzLabelSegment[above,sloped , font=\scriptsize](F,E){\SI{6}{\centi\meter}}
		\tkzLabelSegment[above,sloped , font=\scriptsize](D,F){$y$}
	\end{scope}
\end{tikzpicture}  
}  \parbox{0.55\linewidth}{  
\begin{enumerate}[leftmargin=*, label=\alph*)]
	\item Justifier que les triangles $ACB$ et $FED$ sont semblables.
	\item Écrire les égalités des rapports entre les côtés homologues.
	\item Calculer les longueurs $x$ et $y$.
\end{enumerate}
} 
\end{exercise}
\vfill 
\begin{exercise}\hfill 
	
 \noindent\parbox{0.4\linewidth}{\centering
 	\begin{tikzpicture}[scale=1.0] 
 		\tkzSetUpPoint[size=3, fill=black]
 		\tkzfctset{tan style/.style={->,>=latex, color=red, line width=1.2pt}}   
 		\tkzSetUpLine[line width= 1.2pt, add=.5 and .5] 
 		
 		\tkzDefPoint(0,0){A} 
 		\tkzDefPoint(4,0){B} 
 		\tkzDefPoint(4,3){C} 
 		\tkzDefPoint(0,3){D} 
 		\tkzDrawPolygon[line width=1.2pt](A,B,C,D) 
 		\tkzDefPointBy[projection = onto B--D](C)  \tkzGetPoint{H}
 		\tkzDrawSegments(B,D H,C) 
 		 \tkzDrawSegment[dim={\scriptsize $20$,+10pt,}](D,C)
 		 \tkzDrawSegment[dim={\scriptsize $15$,-10pt,}](B,C)
 		\tkzLabelPoints[above right,font=\scriptsize](C)
 		\tkzLabelPoints[below,font=\scriptsize](B)
 		\tkzLabelPoints[left,font=\scriptsize](A)
 		\tkzLabelPoints[left,font=\scriptsize](D)
 		\tkzLabelPoints[below,font=\scriptsize](H)
 		\tkzLabelSegment[above, sloped , font=\scriptsize](B,H){$c$} 
 		\tkzLabelSegment[above, sloped , font=\scriptsize](H,C){$b$}
 		\tkzLabelSegment[above,sloped ,  font=\scriptsize](D,H){$a$}
 		\tkzMarkRightAngle[size=0.3](C,H,D)
 	\end{tikzpicture}  
 } \parbox{0.55\linewidth}{  
	\begin{enumerate}[leftmargin=*, label=\alph*)]
		\item Calculer la longueur de la diagonale du rectangle $ABCD$.
		\item Justifier que les triangles $ABD$ et $DHC$ sont semblables.
		\item Écrire les égalités des côtés homologues.
		\item Déduire les valeurs de $a$, $b$ et $c$.
	\end{enumerate}
} 
\end{exercise}
\vfill 
%----------------------------------------------------------------------------------------
%	 
%----------------------------------------------------------------------------------------

\begin{exercise}[Quadrature du rectangle] \hfill 
	
	\noindent Sur la figure ci-dessous, le triangle $ABC$ est rectangle en $A$, et $(AH)$ est perpendiculaire à $(BC)$. 
	
	\noindent\parbox[position=t]{0.6 \linewidth}{  
		\begin{enumerate} [leftmargin=*, label=\alph*)]
			\item Démontrer que les triangles $ACH$ et $ABH$ sont semblables. 
			\item Écrire les rapports égaux.
			\item En déduire que $x=\sqrt{ab}$. 
		\end{enumerate}  
	}\hfill
	\noindent\parbox[position=t]{0.35 \linewidth}{  \centering
		\begin{tikzpicture}[scale=0.8]
			\tkzSetUpPoint[size=2, fill=black]
			\tkzfctset{tan style/.style={->,>=latex, color=red, line width=1.2pt}}   
			\tkzSetUpLine[line width= 1.2pt, add=.5 and .5]   
			% 
			\tkzDefPoint(0,0){B} 
			\tkzDefPoint(1.5,0){H}
			\tkzDefPoint(7,0){C}  
			\tkzDefPoint(1.5,3.2){A}  
			\tkzDrawPoints(A,B,C, H) 
			\tkzDrawSegments(A,H H,C C,A A,H A,B B,H)
			\tkzLabelPoints[above,font=\scriptsize](A)
			\tkzLabelPoints[below,font=\scriptsize](B,H,C)
			\tkzLabelSegment[right, font=\scriptsize](A,H){  $x$}
			\tkzLabelSegment[below, font=\scriptsize](H,B){  $a$}
			\tkzLabelSegment[below, font=\scriptsize](H,C){  $b$}
			\tkzMarkRightAngle[ size=0.3](A,H,C) 
			\tkzMarkRightAngle[ size=0.3](B,A,C)	
			%
		\end{tikzpicture}
	}
\end{exercise}  
%----------------------------------------------------------------------------------------
%	 
%----------------------------------------------------------------------------------------

\newpage 
%\begin{exercise}[théorème de l'angle inscrit dans un demi-cercle] \label{chapitre02:exercice:calcullittéraletgeometrie01} 
%	Nous voulons démontrer le théorème suivant :
%\begin{tBox}
%	Si le point $A$ appartient au  cercle de diamètre $[BC]$ alors le triangle $ABC$ est rectangle en $A$.
%	\end{tBox} 
% \noindent\parbox{0.55\linewidth}{
% $A$ est sur le cercle de diamètre $[BC]$ : $OA=OB=OC$. 
%
% On pose $x=\widehat{BAO}$ et $y=\widehat{OAC}$.   
%\begin{enumerate}[leftmargin=*,label=\alph*)]
%	\item Exprimer les angles du triangle $OAB$ à l'aide de $x$.
%	\item En déduire que $\widehat{AOC}=2x$.
%	\item Exprimer $\widehat{AOB}$ à l'aide de $y$.
%	\item Montrer que $x+y=\ang{90}$. 
%\end{enumerate}  
%}\parbox{0.4\linewidth}{\centering
%		\begin{tikzpicture}[scale=1.0]
%			\def \a{3.0};
%			\tkzSetUpPoint[size=3, fill=black]
%			\tkzfctset{tan style/.style={->,>=latex, color=red, line width=1.2pt}} 
%			\tkzDefPoint(0,0){B} 
%			\tkzDefPoint(2*\a,0){C}
%			\tkzDefPoint(\a,0){O}  
%			\tkzDefShiftPoint[O](130:\a){A}
%			\tkzDrawSector[fill = white, line width=1.2pt](O,C)(B)
%			\tkzDrawPolygon[ line width=1.0pt](A,B,C)
%			\tkzDrawSegments[ line width=0.75pt](O,A O,B O,C)
%			\tkzMarkSegments[mark=s||,line     width=0.75pt](O,A O,B O,C)
%			\tkzMarkAngle[mark=|,mkcolor=black, thick, size=0.6, fill=red!20, opacity=0.6](B,A,O)
%			\tkzLabelAngle[pos=1.0,circle](B,A,O){$x$}
%			\tkzMarkAngle[mark=,mkcolor=black, thick, size=0.6, fill=blue!20, opacity=0.6](O,A,C)
%			\tkzLabelAngle[pos=1.1,circle](O,A,C){$y$}
%			%\tkzLabelAngle[pos=1.1,circle](A,C,O){$y$}
%			%\tkzLabelAngle[pos=0.5,circle](O,B,A){$x$}
%			
%			%\tkzLabelAngle[pos=0.9,circle](A,O,B){$2y$}
%			%\tkzLabelAngle[pos=0.5,circle](C,O,A){$2x$}
%			
%			\tkzDrawPoints(A,B,C, O) 
%			\tkzLabelPoints[below,font=\scriptsize](C,O,B) 
%			\tkzLabelPoints[above,font=\scriptsize](A)  
%		\end{tikzpicture}   
%} 
%\end{exercise}
\begin{exercise}[théorème de l'angle au centre] \label{chapitre02:exercice:calcullittéraletgeometrie02}  
	Soit un cercle de centre $O$ passant par les points $A$, $B$ et $C$. Nous voulons démontrer le théorème suivant :
	\begin{tBox}
		Dans un cercle, un angle au centre mesure le double d'un angle inscrit interceptant le même arc.
	\end{tBox}
	
	\noindent\parbox{0.6\linewidth}{
		Pour simplifier on suppose que le centre $O$ est intérieur à l'angle aigu $\widehat{BAC}$.
		Les angles $\widehat{BAC}$ et $\widehat{BOC}$ (intérieur) interceptent le même arc de cercle $ {BC}$.
		
		On pose $x=\widehat{BAO}$ et $y=\widehat{OAC}$.  
		\begin{enumerate}[leftmargin=*,label=\alph*)]
			\item Exprimer les angles du triangle $OAB$ à l'aide de $x$.
			\item Exprimer les angles du triangle $AOC$ à l'aide de $y$. 
			\item Montrer que la mesure de l'angle au centre recherché est égal à $2\widehat{BAC}$. 
		\end{enumerate}  
	}
	\hfill \parbox{0.35\linewidth}{\centering
		\begin{tikzpicture}[scale=1.0]
			\def \a{2.5};
			\tkzSetUpPoint[size=3, fill=black]
			\tkzfctset{tan style/.style={->,>=latex, color=red, line width=1.2pt}} 
			\tkzDefPoint(\a,0){O}  
			\tkzDefShiftPoint[O](-130:\a){B}
			\tkzDefShiftPoint[O](-10:\a){C}  
			\tkzDefShiftPoint[O](120:\a){A}
			\tkzDrawSector[fill = white, line width=1.2pt](O,C)(B)
			%		\tkzDrawPolygon[ line width=1.0pt](A,B,C)
			\tkzDrawArc[line width=1.2pt, dashed](O,B)(C)
			%		\tkzDrawCircle[fill = white, line width=1.2pt](O,A)
			\tkzDrawSegments[ line width=1.0pt](A,B A,C)
			\tkzDrawSegments[ line width=0.75pt](O,A O,B O,C)
			\tkzMarkSegments[mark=s||,line     width=0.75pt](O,A O,B O,C)
			\tkzMarkAngle[mark=none,mkcolor=black, thick, size=0.6, fill=red!20, opacity=0.6](B,A,O)
			\tkzLabelAngle[pos=1.0,circle](B,A,O){$x$}
			\tkzMarkAngle[mark=none,mkcolor=black, thick, size=0.6, fill=blue!20, opacity=0.6](O,A,C)
			\tkzMarkAngle[mark=none,mkcolor=black, thick, size=0.6, fill=blue!20, opacity=0.6](B,O,C)
			\tkzLabelAngle[pos=1.1,circle](B,O,C){$?$}
			\tkzLabelAngle[pos=1.1,circle](O,A,C){$y$}
			%\tkzLabelAngle[pos=1.1,circle](A,C,O){$y$}
			%\tkzLabelAngle[pos=0.5,circle](O,B,A){$x$}
			
			%\tkzLabelAngle[pos=0.9,circle](A,O,B){$2y$}
			%\tkzLabelAngle[pos=0.5,circle](C,O,A){$2x$}
			
			\tkzDrawPoints(A,B,C, O) 
			\tkzLabelPoints[below,font=\scriptsize]( O,B) 
			\tkzLabelPoints[above,font=\scriptsize](A) 
			\tkzLabelPoints[right,font=\scriptsize](C)   
		\end{tikzpicture}  
	} 
\end{exercise}
\begin{theorem}[théorème de l'angle inscrit dans un demi-cercle]
	Si le point $A$ appartient au  cercle de diamètre $[BC]$ alors le triangle $ABC$ est rectangle en $A$.
\end{theorem} 
\noindent
\begin{tikzpicture}[scale=1.0]
	\def \a{2.4};
	\tkzSetUpPoint[size=3, fill=black]
	\tkzfctset{tan style/.style={->,>=latex, color=red, line width=1.2pt}} 
	\tkzDefPoint(0,0){B} 
	\tkzDefPoint(2*\a,0){C}
	\tkzDefPoint(\a,0){O}  
	\tkzDefShiftPoint[O](130:\a){A}
	\tkzDrawSector[fill = white, line width=1.2pt](O,C)(B)
	\tkzDrawPolygon[ line width=1.0pt](A,B,C)
	\tkzDrawSegments[ line width=0.75pt](O,B O,C)
	%		\tkzMarkSegments[mark=s||,line     width=0.75pt](O,A O,B O,C)
	\tkzMarkAngle[<->,>=latex,mkcolor=black, thick, size=0.6, fill=red!20, opacity=0.6](B,O,C)
	\tkzLabelAngle[pos=1.0,circle](B,O,C){\ang{180}}
	\tkzMarkRightAngle[size=0.4](B,A,C)
	%		\tkzMarkAngle[mark=|,mkcolor=black, thick, size=0.6, fill=red!20, opacity=0.6](B,A,O)
	%		\tkzLabelAngle[pos=1.0,circle](B,A,O){$x$}
	%		\tkzMarkAngle[mark=,mkcolor=black, thick, size=0.6, fill=blue!20, opacity=0.6](O,A,C)
	%		\tkzLabelAngle[pos=1.1,circle](O,A,C){$y$}
	%\tkzLabelAngle[pos=1.1,circle](A,C,O){$y$}
	%\tkzLabelAngle[pos=0.5,circle](O,B,A){$x$}
	
	%\tkzLabelAngle[pos=0.9,circle](A,O,B){$2y$}
	%\tkzLabelAngle[pos=0.5,circle](C,O,A){$2x$}
	
	\tkzDrawPoints(A,B,C, O) 
	\tkzLabelPoints[below,font=\scriptsize](C,O,B) 
	\tkzLabelPoints[above,font=\scriptsize](A)  
\end{tikzpicture} \hfill 
\begin{tikzpicture}[scale=1.0]
	\def \a{2.4};
	\tkzSetUpPoint[size=3, fill=black]
	\tkzfctset{tan style/.style={->,>=latex, color=red, line width=1.2pt}} 
	\tkzDefPoint(\a,0){O}  
	\tkzDefShiftPoint[O](-130:\a){B}
	\tkzDefShiftPoint[O](-30:\a){C}  
	\tkzDefShiftPoint[O](120:\a){A}
	\tkzDefShiftPoint[O](10:\a){D}
	\tkzDefShiftPoint[O](65:\a){E}
	\tkzDrawArc[fill = white, line width=1.2pt](O,C)(B)
	%		\tkzDrawPolygon[ line width=1.0pt](A,B,C)
	\tkzDrawArc[line width=1.2pt, dashed](O,B)(C)
	%		\tkzDrawCircle[fill = white, line width=1.2pt](O,A)
	\tkzDrawSegments[ line width=1.0pt](A,B A,C)
	\tkzDrawSegments[ line width=1.0pt](D,B D,C)
	\tkzDrawSegments[ line width=1.0pt](E,B E,C)
	\tkzDrawSegments[ line width=1.0pt,dashed](O,B O,C)
	%	\tkzDrawSegments[ line width=0.75pt](O,A O,B O,C) 
	\tkzMarkAngle[mark=|,mkcolor=black, thick, size=0.6, fill=red!20, opacity=0.6](B,A,C) 
	\tkzMarkAngle[mark=|,mkcolor=black, thick, size=0.6, fill=red!20, opacity=0.6](B,D,C) 
	\tkzMarkAngle[mark=|,mkcolor=black, thick, size=0.6, fill=red!20, opacity=0.6](B,E,C) 
	\tkzMarkAngle[mark=none,mkcolor=black, thick, size=0.6, fill=blue!20, opacity=0.6](B,O,C)
	%	\tkzLabelAngle[pos=1.0,circle](B,A,O){$x$}
	%	\tkzLabelAngle[pos=1.1,circle](B,O,C){$?$}
	%	\tkzLabelAngle[pos=1.1,circle](O,A,C){$y$}
	%\tkzLabelAngle[pos=1.1,circle](A,C,O){$y$}
	%\tkzLabelAngle[pos=0.5,circle](O,B,A){$x$}
	
	%\tkzLabelAngle[pos=0.9,circle](A,O,B){$2y$}
	%\tkzLabelAngle[pos=0.5,circle](C,O,A){$2x$}
	
	\tkzDrawPoints(A,B,C, O) 
	\tkzLabelPoints[below,font=\scriptsize]( O,B) 
	\tkzLabelPoints[above,font=\scriptsize](A) 
	\tkzLabelPoints[right,font=\scriptsize](C)   
\end{tikzpicture}    
\begin{tikzpicture}[scale=1.0]
	\def \a{2.4};
	\tkzSetUpPoint[size=3, fill=black]
	\tkzfctset{tan style/.style={->,>=latex, color=red, line width=1.2pt}} 
	\tkzDefPoint(\a,0){O}  
	\tkzDefShiftPoint[O](-130:\a){B}
	\tkzDefShiftPoint[O](-30:\a){C}  
	\tkzDefShiftPoint[O](-92:\a){A}
	\tkzDefShiftPoint[O](-55:\a){D}
	\tkzDefShiftPoint[O](65:\a){E}
	\tkzDrawArc[fill = white, line width=1.2pt, dashed](O,C)(B)
	%		\tkzDrawPolygon[ line width=1.0pt](A,B,C)
	\tkzDrawArc[line width=1.2pt](O,B)(C)
	%		\tkzDrawCircle[fill = white, line width=1.2pt](O,A)
	\tkzDrawSegments[ line width=1.0pt](A,B A,C)
	\tkzDrawSegments[ line width=1.0pt](D,B D,C)
	%	\tkzDrawSegments[ line width=1.0pt](E,B E,C)
	\tkzDrawSegments[ line width=1.0pt,dashed](O,B O,C)
	%	\tkzDrawSegments[ line width=0.75pt](O,A O,B O,C) 
	\tkzMarkAngle[mark=|,mkcolor=black, thick, size=0.3, fill=red!20, opacity=0.6](C,A,B) 
	\tkzMarkAngle[mark=|,mkcolor=black, thick, size=0.3, fill=red!20, opacity=0.6](C,D,B) 
	%	\tkzMarkAngle[mark=|,mkcolor=black, thick, size=0.6, fill=red!20, opacity=0.6](B,E,C) 
	\tkzMarkAngle[mark=none,mkcolor=black, thick, size=0.6, fill=blue!20, opacity=0.6](C,O,B)
	%	\tkzLabelAngle[pos=1.0,circle](B,A,O){$x$}
	%	\tkzLabelAngle[pos=1.1,circle](B,O,C){$?$}
	%	\tkzLabelAngle[pos=1.1,circle](O,A,C){$y$}
	%\tkzLabelAngle[pos=1.1,circle](A,C,O){$y$}
	%\tkzLabelAngle[pos=0.5,circle](O,B,A){$x$}
	
	%\tkzLabelAngle[pos=0.9,circle](A,O,B){$2y$}
	%\tkzLabelAngle[pos=0.5,circle](C,O,A){$2x$}
	
	\tkzDrawPoints(A,B,C, O) 
	\tkzLabelPoints[below,font=\scriptsize]( O,B) 
	\tkzLabelPoints[below,font=\scriptsize](A) 
	\tkzLabelPoints[right,font=\scriptsize](C)   
\end{tikzpicture}   



\begin{theorem}[Théorème de l'angle inscrit]
	Les angles inscrits interceptant le même arc de cercle ont la même mesure
\end{theorem}
\begin{exercise}\hfill 
	
\noindent\parbox{0.4\linewidth}{
\begin{tikzpicture}[scale=1]
	\def \a{2.4};
	\tkzSetUpPoint[size=3, fill=black]
	\tkzfctset{tan style/.style={->,>=latex, color=red, line width=1.2pt}} 
	\tkzDefPoint(\a,0){O} 
	\tkzDrawPoint(O) 
	
	\tkzDefCircle[R](O,\a) \tkzGetPoint{a} 
	\tkzDrawCircle[line width=1.2pt](O,a)
	
	\tkzDefShiftPoint[O](-130:\a){B}
	\tkzDefShiftPoint[O](120:\a){C}  
	\tkzDefShiftPoint[O](26:\a){A}
	\tkzDefShiftPoint[O](-35:\a){D}
	\tkzDrawSegments(A,B C,D)
	\tkzDrawSegments[dashed](C,B A,D)
	\tkzInterLL(A,B)(C,D)\tkzGetPoint{E}
	
	\tkzLabelSegment[above](A,E){\scriptsize$a$}
	\tkzLabelSegment[above](B,E){\scriptsize$b$}
	\tkzLabelSegment[right](C,E){\scriptsize$c$}
	\tkzLabelSegment[right](D,E){\scriptsize$d$}  
	
	\tkzLabelPoints[left,font=\scriptsize](O) 
	\tkzLabelPoints[right,font=\scriptsize](A) 
	\tkzLabelPoints[below left,font=\scriptsize](B) 
	\tkzLabelPoints[above left,font=\scriptsize](C)  
	\tkzLabelPoints[below right,font=\scriptsize](D)   
	\tkzLabelPoints[above,font=\scriptsize](E) 
%	\tkzDefShiftPoint[O](65:\a){E}
%	\tkzDrawArc[fill = white, line width=1.2pt](O,C)(B)
%	%		\tkzDrawPolygon[ line width=1.0pt](A,B,C)
%	\tkzDrawArc[line width=1.2pt, dashed](O,B)(C)
%	%		\tkzDrawCircle[fill = white, line width=1.2pt](O,A)
%	\tkzDrawSegments[ line width=1.0pt](A,B A,C)
%	\tkzDrawSegments[ line width=1.0pt](D,B D,C)
%	\tkzDrawSegments[ line width=1.0pt](E,B E,C)
%	\tkzDrawSegments[ line width=1.0pt,dashed](O,B O,C)
%	%	\tkzDrawSegments[ line width=0.75pt](O,A O,B O,C) 
%	\tkzMarkAngle[mark=|,mkcolor=black, thick, size=0.6, fill=red!20, opacity=0.6](B,A,C) 
%	\tkzMarkAngle[mark=|,mkcolor=black, thick, size=0.6, fill=red!20, opacity=0.6](B,D,C) 
%	\tkzMarkAngle[mark=|,mkcolor=black, thick, size=0.6, fill=red!20, opacity=0.6](B,E,C) 
%	\tkzMarkAngle[mark=none,mkcolor=black, thick, size=0.6, fill=blue!20, opacity=0.6](B,O,C)
%	%	\tkzLabelAngle[pos=1.0,circle](B,A,O){$x$}
%	%	\tkzLabelAngle[pos=1.1,circle](B,O,C){$?$}
%	%	\tkzLabelAngle[pos=1.1,circle](O,A,C){$y$}
%	%\tkzLabelAngle[pos=1.1,circle](A,C,O){$y$}
%	%\tkzLabelAngle[pos=0.5,circle](O,B,A){$x$}
%	
%	%\tkzLabelAngle[pos=0.9,circle](A,O,B){$2y$}
%	%\tkzLabelAngle[pos=0.5,circle](C,O,A){$2x$}
%	
%	\tkzDrawPoints(A,B,C, O)   
\end{tikzpicture}   
}\hfill
\parbox{0.6\linewidth}{
$A$, $B$, $C$ et $D$ sont des points d'un cercle de centre $O$. On suppose que les cordes $[AB]$ et $[CD]$ se coupent en $E$ situé à l'intérieur du cercle.
\begin{enumerate}[label=\alph*), leftmargin=*]
	\item À l'aide du théorème de l'angle inscrit, justifier que $\widehat{EDA}=\widehat{EBD}$.
	\item Montrer que les triangles $EAD$ et $EBC$ sont semblables.
	\item Écrire les égalités des rapports des côtés homologues.
	\item En déduire que $ab=cd$.
\end{enumerate}
}
\end{exercise}

%----------------------------------------------------------------------------------------
%	 
%----------------------------------------------------------------------------------------
 

%----------------------------------------------------------------------------------------
%	 
%----------------------------------------------------------------------------------------
\newpage 
\loadgeometry{marginlayout}  
\pagestyle{scrheadings} % Print headers again  
\KOMAoptions{
	headwidth=textwithmarginpar,
	headsepline=true,
	headsepline= 0.5pt,
	footwidth=textwithmarginpar,
}   
\onehalfspacing  
\section{Triangles rectangles}
\begin{definition}
	L'hypoténuse d'un triangle rectangle est le côté opposé à  l'angle droit.
\end{definition}
\begin{marginfigure}[-1cm]
	\begin{tikzpicture}[scale=1.0]
		\tkzDefPoint(0,0){A} \tkzDefPoint(4,0){B}
		\tkzDefTriangle[pythagore,swap](A,B)\tkzGetPoint{C}
		\tkzDrawPolygon(A,B,C) \tkzDrawPoints(A,B,C)
		\tkzLabelPoints[below left](A)
		\tkzLabelPoints[below right](B)
		\tkzLabelPoints[above](C)
		%		\tkzMarkAngle[ arc=l, size=0.5](B,A,C)
		\tkzMarkRightAngle[ pos=.15, size=0.4](C,B,A)
		%		\tkzLabelAngle[pos=1.5,circle](B,A,C){$\alpha$}
		%		\tkzLabelSegment[below](A,B){\small adjacent à $\alpha$}
		\tkzLabelSegment[ above , sloped](A,C){hypoténuse}
		%		\tkzLabelSegment[below, sloped](B,C){opposé à $\alpha$}
	\end{tikzpicture} 
\end{marginfigure}
\begin{theorem}[Théorème de Pythagore]  
	Dans un triangle rectangle, le carré de l'hypoténuse est égal à  la somme des carrés des côtés de l'angle droit.
\end{theorem}
\begin{remark} Conséquence : l'hypoténuse est bien le plus grand côté.
	\textcolor{white}{
		\begin{align*}
			AB^2 + BC^2 & = AC^2 && \text{$AB$, $BC$ et $AC\geqslant0$}\\
			AB^2 & \leqslant AC^2 \\
			AB & \leqslant AC
		\end{align*}
	}
\end{remark}

%----------------------------------------------------------------------------------------
%	 
%----------------------------------------------------------------------------------------
%\newpage 
\begin{marginfigure} 
	\begin{tikzpicture}[scale=1]
		\tkzSetUpPoint[size=3, fill=black]
		\tkzfctset{tan style/.style={->,>=latex, color=red, line width=1.2pt}}   
		\tkzSetUpLine[line width= 1.2pt, add=.5 and .5] 
		%		 
		\tkzDefPoint(0,0){B}  
		\tkzDefShiftPoint[B](40:3.0){A}
		\tkzDefShiftPoint[B](15:4.0){C}
		\tkzDrawLine[add=0.2 and 0.2](B,C)
		\tkzDefLine[orthogonal=through A](B,C) \tkzGetPoint{a}
		\tkzLabelLine[pos=1.25,above left](B,C){$(\Delta)$}
		\tkzInterLL(B,C)(A,a) \tkzGetPoint{H}
		\tkzDrawLine[dashed](A,H)
		\tkzMarkRightAngles[size=0.2](C,H,A)
		\tkzDrawPoints(A,H)
		\tkzLabelPoints[below left,font=\scriptsize](A)  
		\tkzLabelPoints[below right,font=\scriptsize](H)   
	\end{tikzpicture} 
\end{marginfigure}  
\begin{definition} Soit une droite $(\Delta)$ et un point $A$ du plan. Le \textbf{projeté orthogonal} $H$ de $A$ sur $(\Delta)$ est le point d'intersection de $\Delta$ et de la perpendiculaire à $(\Delta)$ passant par $A$.
\end{definition} 
\begin{theorem}
	Soit une droite $(\Delta)$ et $A$ un point n'appartenant pas à $(\Delta)$. $H$ le projeté orthogonal de $A$ sur $(\Delta)$.\\
	La distance ente $A$ et la droite $(\Delta)$ est égale à $AH$. C'est la plus petite distance entre $A$ et un point de la droite $(\Delta)$. 
\end{theorem} 
\begin{proof} \emph{au programme} ~ Pour tout point $M\in(Delta)$, $AMH$ est un triangle rectangle d'hypoténuse $AM$, et $AM\geqslant AH$.\hfill 
	\begin{marginfigure}[-1cm] 
		\captionsetup{type=figure} 
		\begin{tikzpicture}[scale=1]
			\tkzSetUpPoint[size=3, fill=black]
			\tkzfctset{tan style/.style={->,>=latex, color=red, line width=1.2pt}}   
			\tkzSetUpLine[line width= 1.2pt, add=.5 and .5] 
			%		 
			\tkzDefPoint(0,0){M}  
			\tkzDefShiftPoint[M](45:3.0){A}
			\tkzDefShiftPoint[M](15:4.0){C}
			\tkzDrawLine[add=0.2 and 0.2](B,C)
			\tkzDefLine[orthogonal=through A](M,C) \tkzGetPoint{a}
			\tkzLabelLine[pos=1.25,above left](M,C){$(\Delta)$}
			\tkzInterLL(M,C)(A,a) \tkzGetPoint{H}
			\tkzDrawLine[dashed](A,H)
			\tkzMarkRightAngles[size=0.2](C,H,A)
			\tkzDrawPoints(A,H,M)
			\tkzLabelPoints[above right,font=\scriptsize](A)  
			\tkzLabelPoints[below right,font=\scriptsize](H)   
			\tkzLabelPoints[below,font=\scriptsize](M)  
			\tkzDrawSegments[dashed](A,M)
			\tkzLabelSegment[font=\scriptsize](A,H){$\delta$}
		\end{tikzpicture} 
	\end{marginfigure}
	
	
%	\vfill 
	
\end{proof}
% \begin{remark} On note que le cercle de centre $A$ et de rayon $AH$ touche la droite $(\Delta)$ en un unique point $H$. On dit que la droite $(\Delta)$ est tangente au cercle.
% \end{remark}
% \begin{definition}	Soit un cercle de centre $O$ et de rayon $R$. En chacun de ses points, le cercle admet une tangente. 
% 	
% 	La tangente en $M$ au cercle est la droite passant par $M$ et perpendiculaire au rayon $OM$.
% 	
% 	La distance entre une tangente et le centre du cercle est égale au rayon $R$.
% \end{definition}
\begin{figure}[hb]\centering 
	\begin{tikzpicture}[scale=1]
		\tkzSetUpPoint[size=3, fill=black]
		\tkzfctset{tan style/.style={->,>=latex, color=red, line width=1.2pt}}   
		\tkzSetUpLine[line width= 1.2pt, add=.5 and .5] 
		%		 
		\tkzDefPoint(0,0){O}
		\tkzDefPoint(3,3){E}
		\tkzDefShiftPoint[O](13:2.5){M}
		%		\tkzDefRandPointOn[circle=center O radius 2.5cm]
		%		\tkzGetPoint{M}
		\tkzDefShiftPoint[O](150:2.5){N}
		%		\tkzDefRandPointOn[circle=center O radius 2.5cm]
		%		\tkzGetPoint{N}
		
		\tkzDefShiftPoint[O](-50:2.5){A}
		\tkzDefShiftPoint[O](-120:2.5){B}
		\tkzDrawSegments(O,M O,N)
		\tkzDrawCircle(O,M) 
		\tkzDefTangent[at=M](O)
		\tkzGetPoint{h}
		\tkzDefTangent[at=N](O)
		\tkzGetPoint{i}
		\tkzDrawLine[add = 2 and 1.5](M,h) 
		\tkzLabelLine[pos=1.25,above left](M,h){$(\Delta)$}
		\tkzDrawLine[add = 2 and 2](N,i)
		\tkzLabelLine[pos=1.25,above left](N,i){$(\Delta')$}
		\tkzDrawLine[dashed, add = 0.5 and 0.5](A,B)
		\tkzMarkRightAngle[](O,M,h)
		\tkzMarkRightAngle[](O,N,i)
		\tkzLabelPoints[above right,font=\scriptsize](O)  
		\tkzLabelPoints[below right,font=\scriptsize](M)  
		\tkzLabelPoints[above left,font=\scriptsize](N)  
		\tkzLabelPoints[below left,font=\scriptsize](B)  
		\tkzLabelPoints[below right,font=\scriptsize](A)  
		\tkzLabelSegments[font=\scriptsize](O,N O,M){$R$}    
		\tkzLabelCircle[font=\scriptsize,  above](O,M)(80){$\mathscr{C}$}
	\end{tikzpicture} 
	\caption{La droite $(AB)$ est sécante au cercle $\mathscr{C}$.\\
		Les droites $(\Delta)$ et $(\Delta')$ sont dites tangentes au cercle. \\
		La tangente en $M$ au cercle en $M$ est perpendiculaire au rayon $OM$.\\
		La distance entre une tangente et le centre du cercle est égale au rayon $R$.\\
		$M$ est le projeté orthogonal du centre $O$ sur droite $(\Delta)$.
	} 
\end{figure}

%----------------------------------------------------------------------------------------
%	
%----------------------------------------------------------------------------------------
\newpage 



\begin{definition} Dans un triangle rectangle d'hypoténuse $1$. $\alpha$ un angle aigu.  On designe par :
	\begin{enumerate}[label=\textbullet, leftmargin=*]
		\item $\cos(\alpha)$ la longueur du côté adjacent à $\alpha$
		\item $\sin(\alpha)$ la longueur du côté opposé à $\alpha$
	\end{enumerate}
	
	\begin{marginfigure} 
	\begin{tikzpicture}[scale=1.0]
		\tkzDefPoint(0,0){A} \tkzDefPoint(4,0){B}
		\tkzDefTriangle[school](A,B)\tkzGetPoint{C}
		\tkzDrawPolygon(A,B,C) \tkzDrawPoints(A,B,C)
		\tkzLabelPoints[below left](A)
		\tkzLabelPoints[below right](B)
		\tkzLabelPoints[above](C)
		\tkzMarkAngle[mark=none, arc=l, size=0.5](B,A,C)
		\tkzMarkRightAngle[ pos=.0.15, size=0.4](C,B,A)
		\tkzLabelAngle[pos=1.5,circle](B,A,C){$\alpha$}
		\tkzLabelSegment[below](A,B){\small  $\cos(\alpha)$}
		\tkzLabelSegment[ above , sloped](A,C){$1$}
		\tkzLabelSegment[below, sloped](B,C){\small  $\sin(\alpha)$}
	\end{tikzpicture}  
\end{marginfigure} 
\end{definition}    
\begin{theorem} Pour tout valeur de l'angle $\alpha$ on a :
\[
 \cos^2(\alpha) + \sin^2(\alpha) = 1
\] 
\end{theorem}
\begin{proof}[démonstration] (\emph{au programme}) 
	Conséquence directe du théorème de Pythagore ! 
\end{proof}

\begin{figure}[!h]
	\begin{tikzpicture}[scale=3.0]
		\tkzSetUpPoint[size=3, fill=black]
		\tkzfctset{tan style/.style={->,>=latex, color=red, line width=1.2pt}}   
		\tkzSetUpLine[line width= 1.2pt, add=.5 and .5] 
		
		\tkzInit[xmin=-1,xmax=1.0,ymin=-0.1,ymax=1,xstep=1, ystep=1];
		%	\tkzGrid[ color=gray ](-3.0,-2.0)(6.0,5.0);   
		\pgfmathsetmacro{\a}{55/360*2*pi*1}  
		

\tkzDefPoint(0,0){O} 
\tkzDefPoint(1,0){I} 
\tkzDefPoint(-1,0){I'} 
\tkzDefPoint(0,1){J} 
\tkzDefPoint({cos(\a)},{sin(\a)}){A} 
\tkzDrawSegments[](O,A) 
\tkzLabelSegment[above left](O,A){$1$}
\tkzLabelPoints[below left](O){\scriptsize O}
\tkzLabelPoint[  right](A){\scriptsize$A(\cos(\alpha);\sin(\alpha))$}   
\tkzLabelAngle[pos=0.35](I,O,A){$\alpha$} 
\tkzMarkAngle[ arc=l,arrows=->,>=latex, size=0.2cm,mkcolor=red, line width =1.2pt](I,O,A)  
\tkzDefPointWith[orthogonal normed,K=1](O,A)
\tkzGetPoint{N} 
%\tkzDrawCircle[ line width=1.2pt](O,I) 
\tkzDrawArc[ line width=1.2pt](O,I)(I') 
\tkzPointShowCoord(A)%[ xlabel=$\cos(\alpha)$,ylabel=$\sin(\alpha)$, xstyle={below=7pt},ystyle={left=6pt}](A) 
\tkzDrawX[above=5pt, label=$x$, font=\scriptsize, line width=1.2pt, right space=0.2pt] 
\tkzDrawY[right=5pt, label=$y$, font=\scriptsize, line width=1.2pt , up space=0.2pt,]
\end{tikzpicture} 
	\caption{Pour des angles obtus.}
\end{figure}
\begin{marginfigure}[2cm]
	\begin{tikzpicture}[scale=1.0]
		\tkzDefPoint(0,0){A} \tkzDefPoint(4,0){B}
		\tkzDefTriangle[school](A,B)\tkzGetPoint{C}
		\tkzDrawPolygon(A,B,C) \tkzDrawPoints(A,B,C)
		\tkzLabelPoints[below left](A)
		\tkzLabelPoints[below right](B)
		\tkzLabelPoints[above](C)
		\tkzMarkAngle[ mark=none,arc=l, size=0.5](B,A,C)
		\tkzMarkRightAngle[ pos=.15, size=0.4](C,B,A)
		\tkzLabelAngle[pos=1.5,circle](B,A,C){$\alpha$}
		\tkzLabelSegment[below](A,B){\small adjacent à  $\alpha$}
		\tkzLabelSegment[ above , sloped](A,C){hypoténuse}
		\tkzLabelSegment[below, sloped](B,C){opposé à  $\alpha$}
	\end{tikzpicture}  
	\caption{En 2\up{nd}, on calcule des rapports trigonométriques d'angles aigus.}
\end{marginfigure}
\begin{definition}
	Soit le triangle $ABC$ rectangle en $B$. 
	\begin{description} 
		\item[le \textbf{sinus}]   de l'angle $\alpha$   
		$$ \sin(\alpha) = \frac{\text{côté opposé à  $\alpha$}}{\text{hypoténuse}} = \frac{BC}{AC}  \leqslant  1$$ 
		\item[le \textbf{cosinus}] de l'angle $\alpha$ : 
		$$ \cos(\alpha) = \frac{\text{côté adjacent à  $\alpha$}}{\text{hypoténuse}} = \frac{AB}{AC} \leqslant  1 $$ 
		\item[la \textbf{tangente}] de l'angle $\alpha$ :
		$$ \tan(\alpha) = \frac{\text{côté opposé à  $\widehat{BAC}$}}{\text{côté adjacent à  $\alpha$}} = \frac{BC}{AB} $$
	\end{description} 
\end{definition}
 
		\vfill 
		
		
%%----------------------------------------------------------------------------------------
%%	EXERCICES
%%----------------------------------------------------------------------------------------
%
% 
%
%\newpage 
%\loadgeometry{widelayout}  
%\pagestyle{scrheadings} % Print headers again  
%\KOMAoptions{
%	headwidth=textwithmarginpar,
%	headsepline=true,
%	headsepline= 0.5pt,
%	footwidth=textwithmarginpar,
%} 
%\onehalfspacing
%\section{Exercices Triplets pythagoriciens} 
%\begin{exercise}\hfill 
%	\begin{enumerate}[label=\alph*), leftmargin=*]
%		\item \textbf{(Un classique de 3\ieme)}	Montrer que la différence des carrés de deux nombres entiers consécutifs est toujours un nombre impair.
%		\item Choisir un nombre $b$ impair strictement supérieur à 3, et l'élever au carré.
%		\item Trouver 2 nombres entiers consécutifs dont la différence des carrés est égale à $b^2$. 
%	\end{enumerate}
%\end{exercise}
%\begin{tBox}
%	Un triplet Pythagoricien est un triplet de trois nombres \textbf{entiers} $a\leqslant  b<c$ tels que $a^2+b^2=c^2$. 
%\end{tBox} 
% \begin{example}[Nous faisons] Cherchons des triplets pythagoriciens $(a; 15; c)$ :
% 	
% 	\noindent \parbox[t]{0.3\linewidth}{
% 		$\begin{WithArrows}
% 			a^2+ 15^2 &\rule{0pt}{1.8em} = c^2\\
% 			15^2&\rule{0pt}{1.8em} = c^2-a^2\\
% 			225& \rule{0pt}{1.8em} =
% 		\end{WithArrows}$  
% 	}
% 	\parbox[t]{0.65\linewidth}{
% 		$\operatorname{diviseurs(225)}=\{
% 		1; 3; 5; 9; \phantom{15; \qquad 25; \qquad 45 ; \qquad ;75; \qquad ; 225 } \}$  
% 		\begin{description}
% 			\item[$c-a=1$];\rule{0pt}{1.8em} $c+a=225$ :
% 			\item[$c-a=3$];\rule{0pt}{1.8em} $c+a=\qquad$ :
% 			\item[$c-a=5$];\rule{0pt}{1.8em} $c+a=\qquad$ :
% 			\item[$c-a=9$];\rule{0pt}{1.8em} $c+a=\qquad$ :
% 			\item[$c-a=15$];\rule{0pt}{1.8em} $c+a=\qquad$ : 
% 		\end{description}
% 	}
% \end{example} 
%\begin{exercise}[À vous]\label{chapitre05:exercice:tripletspythagoriciens}   Combien pouvez vous trouver de triplets pythagoriciens $(a; 12; c)$ ?
%\end{exercise} 
%%\begin{eBox}
%%	La tablette \href{https://en.wikipedia.org/wiki/Plimpton_322}{Plimpton 322} est la plus connue des  \numprint{500000} tablettes babyloniennes dans la collection G. A. Plimpton à l'université de Columbia. Rédigée dans de l'argile en 1800 av. J.-C., elle comporte un tableau de nombres rangés dans 4 colonnes et 15 lignes. Les valeurs des 3 premières colonnes sont associées à des triplets Pythagoriciens vérifiant $a^2+b^2=c^2$ (la 4\ieme étant le numéro de la ligne).  
%%	
%%	Une \href{https://www.maa.org/sites/default/files/pdf/upload_library/22/Ford/Robson105-120.pdf}{analyse minutieuse} a permis d'identifier la méthode utilisée pour obtenir chacune des lignes. Cet algorithme est abordé dans l'exercice suivant.
%%\end{eBox}
%\begin{exercise}[Comprendre la tablette Plimpton 322]\label{chapitre05:exercice:plimpton322}(figure \ref{chapter05:figure:plimpton322})\hfill  
%	\begin{enumerate}[leftmargin=*,label=\alph*)]
%		\item Montrer que pour toute valeur de $x>0$ le triangle ci-dessous est rectangle.
%		
%		%	\begin{figure*}[hb] 
%		\noindent\parbox{1\linewidth}{\centering
%			\begin{tikzpicture}[scale=.8]
%				%			\tkzInit[xmin=-1.0,ymin=-1.0,xmax=6.5,ymax=2.8]
%				%			\tkzClip
%				\tkzDefPoint(0,0){A} \tkzDefPoint(5,0){C}
%				\tkzDefTriangle[two angles = 27 and 90](A,C)\tkzGetPoint{B}
%				\tkzDrawPolygon[line width=1.2pt](A,B,C) 
%				\tkzLabelSegment[below](A,C){$2$}
%				\tkzLabelSegment[right](B,C){$\abs{x-\frac{1}{x}}$}
%				\tkzLabelSegment[above left](A,B){$x+\frac{1}{x}$} 
%			\end{tikzpicture}   
%		}
%		%	\end{figure*}  
%		\item Expliquer pourquoi cette figure illustre l'inégalité pour tout $x> 0$ on a $2\leqslant x+\frac{1}{x}$.
%		\item Donner, sous forme de fractions irréductibles, les longueurs des côtés du triangle rectangle correspondant à la valeur  $x=\tfrac{12}{5}$.
%		
%		Vérifier qu'il est semblable au triangle de côtés $169\colon 119\colon 120$ (1\iere ligne de la tablette).
%		\item Même question pour $x=\tfrac{9}{4}$  et le triangle $97\colon 65\colon 72$ (11\ieme ligne de la tablette).
%		\item Pouvez vous retrouver les triplets $(a;b;c)$ de la seconde ligne, qui correspond à $x=\tfrac{64}{27}$ ?
%	\end{enumerate} 
%\end{exercise} 
%\begin{exercise}  
%Les formules particulières suivantes génèrent des triplets pythagoriciens $(a;b;c)$ : 
%%\setlength{\abovedisplayskip}{5pt}
%%\setlength{\belowdisplayskip}{5pt}
%%\begin{equation}
%%\label{formules:tripletspythagore}
%
%$
%p;q\in\mathbb{N} \qquad a=p^2-q^2 \qquad b = 2 pq \qquad c = p^2+q^2 
%$
%%\end{equation} 
%	\begin{enumerate}[label=\alph*), leftmargin=*]
%		\item Donner le triplet correspondant à $(p=8;q=7)$ et vérifier l'égalité $a^2+b^2=c^2$ 
%		\item Retrouver le couple d'entiers $(p;q)$ qui donnent le triplet $(a=35;b=12;c=37)$. 
%		\item Développer, réduire et montrer que pour tout $p,q\in\mathbb{N}$ : $a^2+b^2=p^4+2p^2q^2+q^4$.
%		\item  Montrer que pour tout $p,q\in\mathbb{N}$, on a  $a^2+b^2=c^2$.
%	\end{enumerate}
%\end{exercise}
% 
%%----------------------------------------------------------------------------------------
%%	 
%%----------------------------------------------------------------------------------------
%
% \newpage  
%\begin{exercise}\label{chapitre05:exercice:tripletspythagoriciensAlgorithme}\hfill 
%	
%	\begin{enumerate}[label=\arabic*),leftmargin=*]
%		\item Voici 3 scripts Python. Cocher les cases correspondant aux scripts qui affichent le couple indiqué :
%		\begin{multicols}{2} \setlength\columnseprule{.0pt}
%\begin{lstlisting}[ language=python , 
%title=Script A,   
%linewidth=1\columnwidth,
%style=python-code,    
%aboveskip=0.2\topskip,
%belowskip=0.2\topskip,
%%xleftmargin=1.2em,
%lineskip=1.1em,
%tabsize=4, 
%numbers=none,
%%	firstnumber=10  
%]
%for a in range(100) :
%	for b in range(100) :
%		print(a,b) 
%\end{lstlisting}  
%\begin{lstlisting}[ language=python , 
%title=Script B,   
%linewidth=1\columnwidth,
%style=python-code,    
%aboveskip=0.2\topskip,
%belowskip=0.2\topskip,
%%xleftmargin=1.2em,
%lineskip=1.1em,
%tabsize=4, 
%numbers=none,
%%	firstnumber=10  
%] 
%for a in range(100) :
%	for b in range(a,100-a) :
%		print(a,b) 
%		
%\end{lstlisting} 
%
%\begin{lstlisting}[ language=python , 
%title=Script C,   
%linewidth=1\columnwidth,
%style=python-code,    
%aboveskip=0.2\topskip,
%belowskip=0.2\topskip,
%%xleftmargin=1.2em,
%lineskip=1.1em,
%tabsize=4, 
%numbers=none,
%	%	firstnumber=10  
%	] 
%for a in range(100) :
%	b = 100 - a
%	print(a,b)
%\end{lstlisting}
%\begin{lstlisting}[ language=python , 
%title=Script D,   
%linewidth=1\columnwidth,
%style=python-code,    
%aboveskip=0.2\topskip,
%belowskip=0.2\topskip,
%%xleftmargin=1.2em,
%lineskip=1.1em,
%tabsize=4, 
%numbers=none,
%	%	firstnumber=10  
%	] 
%for a in range(100) :
%	for b in range(a,100) :
%		if a**2+b**2 == 10**2 :
%			print(a,b)  
%\end{lstlisting}
%		\end{multicols}
%			 
%\noindent \QCM[Multiple,Alterne,Reponses=4,Stretch=1.1,Largeur=3.0cm,%Solution,%
%Noms={Script A/Script B/Script C/Script D}]{%\textbb{N} marche pas
%{\lstinline|0, 10|} &1&1&0&1, 
%{\lstinline|10, 0|} &1&0&0&0, 
%{\lstinline|1, 99|} &1&0&1&0,
%{\lstinline|3, 7|} &1&1&0&0, 
%{\lstinline|4, 3|} &1&0&0&0, 
%{\lstinline|6, 8|} &1&1&0&1,  
%{\lstinline|8, 6|} &1&0&0&0, 
%%{\lstinline|36, 64|} &1&0&1&0,
%{\lstinline|75, 25|} &1&0&1&0, 
%} 
%\vfill 
%		\item On cherche les triplets pythagoriciens $(a,b,c)$ correspondant à des triangles de périmètre 1000.
% 
%\noindent \begin{lstlisting}[ language=python , 
%%	title=Script D,   
%	linewidth=1\columnwidth,
%	style=python-code,    
%	aboveskip=0.2\topskip,
%	belowskip=0.2\topskip,
%	%xleftmargin=1.2em,
%	lineskip=1.1em,
%	tabsize=4, 
%	numbers=none,
%	%	firstnumber=10  
%	] 
%for a in ................       :
%	for b in ................   :
%		c = ...................
%		if .................... :
%			print(a,b,c)
%\end{lstlisting} 
%\begin{enumerate}[label=\alph*),leftmargin=*]
%	\item Expliquer pourquoi $a$, $b$ et $c$ sont forcement inférieurs à 1000.
%	\item Compléter le script ci-dessous afin qu'il affiche les bons triplets.
%	\item Rentrer le script sur votre pythonette et l'exécuter. \\
%	Vérifier qu'il n'affiche aucune valeur négative. puis relever les triplets solutions.
%\end{enumerate} 
%\end{enumerate}
%\end{exercise}
%%----------------------------------------------------------------------------------------
%%	 
%%----------------------------------------------------------------------------------------
%
%\newpage 
%\begin{proof}[solution de l'exemple et de l'exercice \ref{chapitre05:exercice:tripletspythagoriciens}]\hfill  
%	\begin{enumerate}[label=\alph*),leftmargin=*]
%		\item Pour $b=15$ on trouve $117^2+15^2=118^2$, $36^2+15^2=39^2$, $20^2+15^2=25^2$ et $8^2+15^2=17^2$.
%		\item Pour $b=12$ on trouve $35^2+12^2=37^2$, $16^2+12^2=20^2$, $9^2+12^2=15^2$ et $5^2+12^2=13^2$. 
%	\end{enumerate}
%\end{proof}
% \begin{figure*}[hb]
%	\includegraphics[width=0.60\linewidth]{Plimpton322-Drawing by Eleanor Robson.png} 
%	\caption{La tablette \href{https://en.wikipedia.org/wiki/Plimpton_322}{Plimpton 322} est la plus connue des  \numprint{500000} tablettes babyloniennes dans la collection G. A. Plimpton à l'université de Columbia. Rédigée dans de l'argile en 1800 av. J.-C., elle comporte un tableau de nombres rangés dans 4 colonnes et 15 lignes. Les valeurs des 3 premières colonnes sont associées à des triplets pythagoriciens vérifiant $a^2+b^2=c^2$ (la 4\ieme étant le numéro de la ligne).  
%	Une \href{https://www.maa.org/sites/default/files/pdf/upload_library/22/Ford/Robson105-120.pdf}{analyse minutieuse} a permis d'identifier la méthode utilisée pour obtenir chacune des lignes. Cet algorithme est abordé dans l'exercice \ref{chapitre05:exercice:plimpton322}.}\label{chapter05:figure:plimpton322}
%\end{figure*} 
%\begin{proof}[solution de l'exercice \ref{chapitre05:exercice:tripletspythagoriciensAlgorithme}]\hfill  
%	
%\noindent \QCM[Multiple,Alterne,Reponses=4,Stretch=1.0,Largeur=3.0cm,Solution,%
%Noms={Script A/Script B/Script C/Script D}]{%\textbb{N} marche pas
%	{\lstinline|0, 10|} &1&1&0&1, 
%	{\lstinline|10, 0|} &1&0&0&0, 
%	{\lstinline|1, 99|} &1&0&1&0,
%	{\lstinline|3, 7|} &1&1&0&0, 
%	{\lstinline|4, 3|} &1&0&0&0, 
%	{\lstinline|6, 8|} &1&1&0&1,  
%	{\lstinline|8, 6|} &1&0&0&0, 
%	%{\lstinline|36, 64|} &1&0&1&0,
%	{\lstinline|75, 25|} &1&0&1&0, 
%} 
%
%\begin{lstlisting}[ language=python , 
%	%	title=Script D,   
%	linewidth=1\columnwidth,
%	style=python-code,    
%	aboveskip=0.2\topskip,
%	belowskip=0.2\topskip,
%	%xleftmargin=1.2em,
%%	lineskip=1.0em,
%	tabsize=4, 
%	numbers=none,
%	%	firstnumber=10  
%	] 
%for a in range(1000) :
%	for b in range(a,1000-a)  :
%		c = 1000 - a - b
%		if a**2 + b**2 == c**2 :
%			print(a,b,c)				# programme affiche 200, 375, 475
%\end{lstlisting} 
%\end{proof}


%----------------------------------------------------------------------------------------
%	EXERCICES
%----------------------------------------------------------------------------------------
 
%%----------------------------------------------------------------------------------------
%%	EXERCICES
%%----------------------------------------------------------------------------------------
%
% 
%
\newpage 
\loadgeometry{widelayout}  
\pagestyle{scrheadings} % Print headers again  
\KOMAoptions{
		headwidth=textwithmarginpar,
		headsepline=true,
		headsepline= 0.5pt,
		footwidth=textwithmarginpar,
	} 
\onehalfspacing

\subsection{Exercices : trigonométrie} 
\begin{exercise}[Trigonométrie] \label{chapter05:exercices:trigonometrie:01} Écrire le rapport trigonométrique adapté et calculer les valeurs demandées.
	
	\noindent\parbox[position=t]{0.23\textwidth}{
		%\begin{figure}[H]\centering
		\begin{tikzpicture}[scale=1.0]
			\tkzText(0.5,2.8){$(A)$}
			\tkzSetUpPoint[size=3, fill=black]
			\tkzfctset{tan style/.style={->,>=latex, color=red, line width=1.2pt}}   
			\tkzSetUpLine[line width= 1.2pt, add=.5 and .5] 
			\tkzDefPoint(1,0){A} 
			\tkzDefShiftPoint[A](0:3.3){B}
			\tkzDefTriangle[two angles=90 and 35](A,B)\tkzGetPoint{C}  
			\tkzDrawPolygon[line width=1.2pt](A,B,C)
			%		\tkzLabelPoints[font=\scriptsize](A,B,C)
			\tkzMarkRightAngle[](B,A,C)
			%\tkzMarkAngle[mark=none,-, >=latex, size =0.5](C,B,A)
			%\tkzMarkAngle[mark=none,-, >=latex, size =0.5](A,C,B)
			\tkzLabelAngle[font=\footnotesize,pos=.9](C,B,A){$?$}
			%\tkzLabelAngle[font=\scriptsize,pos=.8](A,C,B){$\ang{55}$}
			%\tkzLabelSegment[above, font=\scriptsize](A,C){$10$} 
			\tkzLabelSegment[below, font=\scriptsize](A,B){\SI{7}{\centi\meter}}
			\tkzLabelSegment[above,sloped,  font=\scriptsize](C,B){\SI{9}{\centi\meter}} 
		\end{tikzpicture}   
	}
%	\hfill 
	\parbox[position=t]{0.23\textwidth}{
		%\begin{figure}[H]\centering
		\begin{tikzpicture}[scale=1.0]
			\tkzText(0.5,2.8){$(B)$}
			\tkzSetUpPoint[size=3, fill=black]
			\tkzfctset{tan style/.style={->,>=latex, color=red, line width=1.2pt}}   
			\tkzSetUpLine[line width= 1.2pt, add=.5 and .5] 
			\tkzDefPoint(1,0){A} 
			\tkzDefShiftPoint[A](0:3.3){B}
			\tkzDefTriangle[two angles=90 and 35](A,B)\tkzGetPoint{C}  
			\tkzDrawPolygon[line width=1.2pt](A,B,C)
			%		\tkzLabelPoints[font=\scriptsize](A,B,C)
			\tkzMarkRightAngle[](B,A,C)
			%\tkzMarkAngle[mark=none,-, >=latex, size =0.5](C,B,A)
			%\tkzMarkAngle[mark=none,-, >=latex, size =0.5](A,C,B)
			%		\tkzLabelAngle[font=\footnotesize,pos=.9](C,B,A){$?$}
			\tkzLabelAngle[font=\footnotesize,pos=.8](A,C,B){$?$}
			%\tkzLabelSegment[above, font=\scriptsize](A,C){$10$} 
			\tkzLabelSegment[below, font=\scriptsize](A,B){\SI{4}{\centi\meter}}
			\tkzLabelSegment[above,sloped,  font=\scriptsize](C,B){\SI{5}{\centi\meter}} 
		\end{tikzpicture}   
	}
%	\hfill 
	\parbox[position=t]{0.23\textwidth}{
		%\begin{figure}[H]\centering
		\begin{tikzpicture}[scale=1.0]
			\tkzText(0.5,2.8){$(C)$}
			\tkzSetUpPoint[size=3, fill=black]
			\tkzfctset{tan style/.style={->,>=latex, color=red, line width=1.2pt}}   
			\tkzSetUpLine[line width= 1.2pt, add=.5 and .5] 
			\tkzDefPoint(1,0){A} 
			\tkzDefShiftPoint[A](0:3.3){B}
			\tkzDefTriangle[two angles=90 and 35](A,B)\tkzGetPoint{C}  
			\tkzDrawPolygon[line width=1.2pt](A,B,C)
			%		\tkzLabelPoints[font=\scriptsize](A,B,C)
			\tkzMarkRightAngle[](B,A,C)
			\tkzMarkAngle[mark=none,-, >=latex, size =0.5](C,B,A)
			%\tkzMarkAngle[mark=none,-, >=latex, size =0.5](A,C,B)
			\tkzLabelAngle[font=\footnotesize,pos=1.0](C,B,A){$\ang{40}$}
			%		\tkzLabelAngle[font=\footnotesize,pos=.8](A,C,B){$40$}
			%\tkzLabelSegment[above, font=\scriptsize](A,C){$10$} 
			\tkzLabelSegment[below, font=\footnotesize](A,B){$?$}
			\tkzLabelSegment[above,sloped,  font=\scriptsize](C,B){\SI{5}{\centi\meter}} 
		\end{tikzpicture}   
	}
%	\hfill 
	\parbox[position=t]{0.23\textwidth}{
		%\begin{figure}[H]\centering
		\begin{tikzpicture}[scale=1.0]
			\tkzText(0.5,2.8){$(D)$}
			\tkzSetUpPoint[size=3, fill=black]
			\tkzfctset{tan style/.style={->,>=latex, color=red, line width=1.2pt}}   
			\tkzSetUpLine[line width= 1.2pt, add=.5 and .5] 
			\tkzDefPoint(1,0){A} 
			\tkzDefShiftPoint[A](0:3.3){B}
			\tkzDefTriangle[two angles=90 and 35](A,B)\tkzGetPoint{C}  
			\tkzDrawPolygon[line width=1.2pt](A,B,C)
			%		\tkzLabelPoints[font=\scriptsize](A,B,C)
			\tkzMarkRightAngle[](B,A,C)
			\tkzMarkAngle[mark=none,-, >=latex, size =0.5](C,B,A)
			%\tkzMarkAngle[mark=none,-, >=latex, size =0.5](A,C,B)
			\tkzLabelAngle[font=\footnotesize,pos=1.0](C,B,A){$\ang{40}$}
			%		\tkzLabelAngle[font=\footnotesize,pos=.8](A,C,B){$40$}
			%\tkzLabelSegment[above, font=\scriptsize](A,C){$10$} 
			\tkzLabelSegment[below, font=\footnotesize](A,B){\SI{5}{\centi\meter}}
			\tkzLabelSegment[above,sloped,  font=\scriptsize](C,B){$?$} 
		\end{tikzpicture}   
	}
	
	\noindent
	\parbox[position=t]{0.23\textwidth}{
		%\begin{figure}[H]\centering
		\begin{tikzpicture}[scale=1.0]
			\tkzText(0.5,2.8){$(E)$}
			\tkzSetUpPoint[size=3, fill=black]
			\tkzfctset{tan style/.style={->,>=latex, color=red, line width=1.2pt}}   
			\tkzSetUpLine[line width= 1.2pt, add=.5 and .5] 
			\tkzDefPoint(1,0){A} 
			\tkzDefShiftPoint[A](0:3.3){B}
			\tkzDefTriangle[two angles=90 and 35](A,B)\tkzGetPoint{C}  
			\tkzDrawPolygon[line width=1.2pt](A,B,C)
			%		\tkzLabelPoints[font=\scriptsize](A,B,C)
			\tkzMarkRightAngle[](B,A,C)
			%		\tkzMarkAngle[mark=none,-, >=latex, size =0.5](C,B,A)
			\tkzMarkAngle[mark=none,-, >=latex, size =0.5](A,C,B)
			%		\tkzLabelAngle[font=\footnotesize,pos=1.0](C,B,A){$\ang{40}$}
			\tkzLabelAngle[font=\footnotesize,pos=.8](A,C,B){$\ang{40}$}
			%\tkzLabelSegment[above, font=\scriptsize](A,C){$10$} 
			\tkzLabelSegment[below, font=\footnotesize](A,B){\SI{3}{\centi\meter}}
			\tkzLabelSegment[above,sloped,  font=\scriptsize](C,B){$?$} 
		\end{tikzpicture}   
	}
%	\hfill
	\parbox[position=t]{0.23\textwidth}{
		%\begin{figure}[H]\centering
		\begin{tikzpicture}[scale=1.0]
			\tkzText(0.5,2.8){$(F)$}
			\tkzSetUpPoint[size=3, fill=black]
			\tkzfctset{tan style/.style={->,>=latex, color=red, line width=1.2pt}}   
			\tkzSetUpLine[line width= 1.2pt, add=.5 and .5] 
			\tkzDefPoint(1,0){A} 
			\tkzDefShiftPoint[A](0:3.3){B}
			\tkzDefTriangle[two angles=90 and 35](A,B)\tkzGetPoint{C}  
			\tkzDrawPolygon[line width=1.2pt](A,B,C)
			%		\tkzLabelPoints[font=\scriptsize](A,B,C)
			\tkzMarkRightAngle[](B,A,C)
			%		\tkzMarkAngle[mark=none,-, >=latex, size =0.5](C,B,A)
			\tkzMarkAngle[mark=none,-, >=latex, size =0.5](A,C,B)
			%		\tkzLabelAngle[font=\footnotesize,pos=1.0](C,B,A){$\ang{40}$}
			\tkzLabelAngle[font=\footnotesize,pos=.8](A,C,B){$?$}
			\tkzLabelSegment[left, font=\scriptsize](A,C){$\SI{3}{\centi\meter}$} 
			\tkzLabelSegment[below, font=\footnotesize](A,B){\SI{4}{\centi\meter}}
			%		\tkzLabelSegment[above,sloped,  font=\scriptsize](C,B){$?$} 
		\end{tikzpicture}   
	} 
%	\hfill
	\parbox[position=t]{0.23\textwidth}{
		%\begin{figure}[H]\centering
		\begin{tikzpicture}[scale=1.0]
			\tkzText(0.5,2.8){$(G)$}
			\tkzSetUpPoint[size=3, fill=black]
			\tkzfctset{tan style/.style={->,>=latex, color=red, line width=1.2pt}}   
			\tkzSetUpLine[line width= 1.2pt, add=.5 and .5] 
			\tkzDefPoint(1,0){A} 
			\tkzDefShiftPoint[A](0:3.3){B}
			\tkzDefTriangle[two angles=35 and 55](A,B)\tkzGetPoint{C}  
			\tkzDrawPolygon[line width=1.2pt](A,B,C)
			%		\tkzLabelPoints[font=\scriptsize](A,B,C)
			%		\tkzMarkRightAngle[](B,A,C)
			\tkzMarkRightAngle[](A,C,B)
			%		\tkzMarkAngle[mark=none,-, >=latex, size =0.5](C,B,A)
			\tkzMarkAngle[mark=none,-, >=latex, size =0.5](B,A,C)
			%		\tkzLabelAngle[font=\footnotesize,pos=1.0](C,B,A){$\ang{40}$}
			%		\tkzLabelAngle[font=\footnotesize,pos=.8](A,C,B){$?$}
			\tkzLabelAngle[font=\footnotesize,pos=.8](B,A,C){$?$}
			\tkzLabelSegment[left, font=\scriptsize](A,C){$\SI{5}{\centi\meter}$} 
			\tkzLabelSegment[below, font=\footnotesize](A,B){\SI{7}{\centi\meter}}
			%		\tkzLabelSegment[above,sloped,  font=\scriptsize](C,B){$?$} 
		\end{tikzpicture}   
	}
%	\hfill
	\parbox[position=t]{0.23\textwidth}{
		%\begin{figure}[H]\centering
		\begin{tikzpicture}[scale=1.0]
			\tkzText(0.5,2.8){$(H)$}
			\tkzSetUpPoint[size=3, fill=black]
			\tkzfctset{tan style/.style={->,>=latex, color=red, line width=1.2pt}}   
			\tkzSetUpLine[line width= 1.2pt, add=.5 and .5] 
			\tkzDefPoint(1,0){A} 
			\tkzDefShiftPoint[A](0:3.3){B}
			\tkzDefTriangle[two angles=35 and 55](A,B)\tkzGetPoint{C}  
			\tkzDrawPolygon[line width=1.2pt](A,B,C)
			%		\tkzLabelPoints[font=\scriptsize](A,B,C)
			%		\tkzMarkRightAngle[](B,A,C)
			\tkzMarkRightAngle[](A,C,B)
			%		\tkzMarkAngle[mark=none,-, >=latex, size =0.5](C,B,A)
			\tkzMarkAngle[mark=none,-, >=latex, size =0.5](B,A,C)
			%		\tkzLabelAngle[font=\footnotesize,pos=1.0](C,B,A){$\ang{40}$}
			%		\tkzLabelAngle[font=\footnotesize,pos=.8](A,C,B){$?$}
			\tkzLabelAngle[font=\footnotesize,pos=.8](B,A,C){$?$}
			%	\tkzLabelSegment[left, font=\scriptsize](A,C){$\SI{5}{\centi\meter}$} 
			\tkzLabelSegment[above right, font=\scriptsize](B,C){$\SI{5}{\centi\meter}$} 
			\tkzLabelSegment[below, font=\footnotesize](A,B){\SI{10}{\centi\meter}}
			%		\tkzLabelSegment[above,sloped,  font=\scriptsize](C,B){$?$} 
		\end{tikzpicture}   
	}
\end{exercise}
\begin{example}[Calcul d'une aire] Calculer l'aire du triangle $ABC$ ci-dessous.
	
	
\noindent\parbox[position=t]{0.3\textwidth}{\centering
	%\begin{figure}[H]\centering
	\begin{tikzpicture}[scale=1.0] 
		\tkzSetUpPoint[size=3, fill=black]
		\tkzfctset{tan style/.style={->,>=latex, color=red, line width=1.2pt}}   
		\tkzSetUpLine[line width= 1.2pt, add=.5 and .5] 
		
		\tkzDefPoint(0,0){A} 
		\tkzDefShiftPoint[A](0:4){B}
		\tkzDefShiftPoint[A](40:3){C}
		% 			\tkzDefTriangle[two angles=90 and 35](A,B)\tkzGetPoint{C}  
		\tkzDrawPolygon[line width=1.2pt](A,B,C)
		%		\tkzLabelPoints[font=\scriptsize](A,B,C)
		% 			\tkzMarkRightAngle[](B,A,C)
		\tkzMarkAngle[mark=none,-, >=latex, size =0.5](B,A,C)
		% 			\tkzMarkAngle[mark=none,-, >=latex, size =0.5](C,B,A)
		%\tkzMarkAngle[mark=none,-, >=latex, size =0.5](A,C,B)
		\tkzLabelAngle[font=\footnotesize,pos=.9](B,A,C){$\ang{40}$}
		% 			\tkzLabelAngle[font=\footnotesize,pos=.9](C,B,A){$?$}
		%\tkzLabelAngle[font=\scriptsize,pos=.8](A,C,B){$\ang{55}$}
		
		\tkzDrawSegment[dim={\scriptsize \SI{6}{\centi\meter},+15pt,}](A,C)
		\tkzDrawSegment[dim={\scriptsize \SI{8}{\centi\meter},-15pt,}](A,B)
		
		
%		\tkzLabelSegment[above, font=\scriptsize,sloped](A,C){\SI{5}{\centi\meter}} 
%		\tkzLabelSegment[below, font=\scriptsize](A,B){\SI{6}{\centi\meter}}

		% 			\tkzLabelSegment[above,sloped,  font=\scriptsize](C,B){\SI{9}{\centi\meter}}  
		\tkzLabelPoints[below left,font=\scriptsize](A)
		\tkzLabelPoints[below right,font=\scriptsize](B)
		\tkzLabelPoints[above,font=\scriptsize](C)  
		\tkzDefPointBy[projection = onto A--B](C) \tkzGetPoint{H}
		\tkzDrawSegment[dashed](C,H)
		\tkzMarkRightAngle[size=0.2](B,H,C)
		\tkzLabelPoints[below left,font=\scriptsize](H)
	\end{tikzpicture}   
} 
\end{example}
\vfill 
\begin{exercise}[À vous]\label{chapitre05:exercice:calculdaires01} Calculer l'aire des triangles suivants :\hfill 
	
\noindent\parbox[position=t]{0.3\textwidth}{\centering
	%\begin{figure}[H]\centering
	\begin{tikzpicture}[scale=1.0] 
		\tkzSetUpPoint[size=3, fill=black]
		\tkzfctset{tan style/.style={->,>=latex, color=red, line width=1.2pt}}   
		\tkzSetUpLine[line width= 1.2pt, add=.5 and .5] 
		
		\tkzDefPoint(0,0){A} 
		\tkzDefShiftPoint[A](0:3){B}
		\tkzDefShiftPoint[A](65:2.5){C}
		% 			\tkzDefTriangle[two angles=90 and 35](A,B)\tkzGetPoint{C}  
		\tkzDrawPolygon[line width=1.2pt](A,B,C)
		%		\tkzLabelPoints[font=\scriptsize](A,B,C)
		% 			\tkzMarkRightAngle[](B,A,C)
		\tkzMarkAngle[mark=none,-, >=latex, size =0.5](B,A,C)
		% 			\tkzMarkAngle[mark=none,-, >=latex, size =0.5](C,B,A)
		%\tkzMarkAngle[mark=none,-, >=latex, size =0.5](A,C,B)
		\tkzLabelAngle[font=\footnotesize,pos=.9](B,A,C){$\ang{65}$}
		% 			\tkzLabelAngle[font=\footnotesize,pos=.9](C,B,A){$?$}
		%\tkzLabelAngle[font=\scriptsize,pos=.8](A,C,B){$\ang{55}$}
		
		\tkzDrawSegment[dim={\scriptsize \SI{5}{\centi\meter},+15pt,}](A,C)
		\tkzDrawSegment[dim={\scriptsize \SI{6}{\centi\meter},-15pt,}](A,B)
		
		
		%		\tkzLabelSegment[above, font=\scriptsize,sloped](A,C){\SI{5}{\centi\meter}} 
		%		\tkzLabelSegment[below, font=\scriptsize](A,B){\SI{6}{\centi\meter}}
		
		% 			\tkzLabelSegment[above,sloped,  font=\scriptsize](C,B){\SI{9}{\centi\meter}}  
		\tkzLabelPoints[below left,font=\scriptsize](A)
		\tkzLabelPoints[below right,font=\scriptsize](B)
		\tkzLabelPoints[above,font=\scriptsize](C)  
		\tkzDefPointBy[projection = onto A--B](C) \tkzGetPoint{H}
		\tkzDrawSegment[dashed](C,H)
		\tkzMarkRightAngle[size=0.2](B,H,C)
		\tkzLabelPoints[below left,font=\scriptsize](H)
	\end{tikzpicture}   
}\hfill \parbox[position=t]{0.3\textwidth}{\centering
%\begin{figure}[H]\centering
\begin{tikzpicture}[scale=1.0] 
	\tkzSetUpPoint[size=3, fill=black]
	\tkzfctset{tan style/.style={->,>=latex, color=red, line width=1.2pt}}   
	\tkzSetUpLine[line width= 1.2pt, add=.5 and .5] 
	
	\tkzDefPoint(0,0){A} 
	\tkzDefShiftPoint[A](-10:5){B}
	\tkzDefShiftPoint[A](28:3){C}
	% 			\tkzDefTriangle[two angles=90 and 35](A,B)\tkzGetPoint{C}  
	\tkzDrawPolygon[line width=1.2pt](A,B,C)
	%		\tkzLabelPoints[font=\scriptsize](A,B,C)
	% 			\tkzMarkRightAngle[](B,A,C)
	\tkzMarkAngle[mark=none,-, >=latex, size =0.5](B,A,C)
	% 			\tkzMarkAngle[mark=none,-, >=latex, size =0.5](C,B,A)
	%\tkzMarkAngle[mark=none,-, >=latex, size =0.5](A,C,B)
	\tkzLabelAngle[font=\footnotesize,pos=.9](B,A,C){$\ang{38}$}
	% 			\tkzLabelAngle[font=\footnotesize,pos=.9](C,B,A){$?$}
	%\tkzLabelAngle[font=\scriptsize,pos=.8](A,C,B){$\ang{55}$}
	
	\tkzDrawSegment[dim={\scriptsize \SI{6}{\centi\meter},+15pt,}](A,C)
	\tkzDrawSegment[dim={\scriptsize \SI{10}{\centi\meter},-15pt,}](A,B)
	
	
	%		\tkzLabelSegment[above, font=\scriptsize,sloped](A,C){\SI{5}{\centi\meter}} 
	%		\tkzLabelSegment[below, font=\scriptsize](A,B){\SI{6}{\centi\meter}}
	
	% 			\tkzLabelSegment[above,sloped,  font=\scriptsize](C,B){\SI{9}{\centi\meter}}  
%	\tkzLabelPoints[below left,font=\scriptsize](A)
%	\tkzLabelPoints[below right,font=\scriptsize](B)
%	\tkzLabelPoints[above,font=\scriptsize](C)  
	\tkzDefPointBy[projection = onto A--B](C) \tkzGetPoint{H}
	\tkzDrawSegment[dashed](C,H)
	\tkzMarkRightAngle[size=0.2](B,H,C)
%	\tkzLabelPoints[below left,font=\scriptsize](H)
\end{tikzpicture}   
}\hfill \parbox[position=t]{0.3\textwidth}{\centering
%\begin{figure}[H]\centering
\begin{tikzpicture}[scale=1.0] 
	\tkzSetUpPoint[size=3, fill=black]
	\tkzfctset{tan style/.style={->,>=latex, color=red, line width=1.2pt}}   
	\tkzSetUpLine[line width= 1.2pt, add=.5 and .5] 
	
	\tkzDefPoint(0,0){A} 
	\tkzDefShiftPoint[A](0:4.5){B}
	\tkzDefShiftPoint[A](-76:3.5){C}
	% 			\tkzDefTriangle[two angles=90 and 35](A,B)\tkzGetPoint{C}  
	\tkzDrawPolygon[line width=1.2pt](A,B,C)
	%		\tkzLabelPoints[font=\scriptsize](A,B,C)
	% 			\tkzMarkRightAngle[](B,A,C)
	\tkzMarkAngle[mark=none,-, >=latex, size =0.5](C,A,B)
	% 			\tkzMarkAngle[mark=none,-, >=latex, size =0.5](C,B,A)
	%\tkzMarkAngle[mark=none,-, >=latex, size =0.5](A,C,B)
	\tkzLabelAngle[font=\footnotesize,pos=.9](C,A,B){$\ang{76}$}
	% 			\tkzLabelAngle[font=\footnotesize,pos=.9](C,B,A){$?$}
	%\tkzLabelAngle[font=\scriptsize,pos=.8](A,C,B){$\ang{55}$}
	
	\tkzDrawSegment[dim={\scriptsize \SI{6}{\centi\meter},-15pt,}](A,C)
	\tkzDrawSegment[dim={\scriptsize \SI{10}{\centi\meter},+15pt,}](A,B)
	
	
	%		\tkzLabelSegment[above, font=\scriptsize,sloped](A,C){\SI{5}{\centi\meter}} 
	%		\tkzLabelSegment[below, font=\scriptsize](A,B){\SI{6}{\centi\meter}}
	
	% 			\tkzLabelSegment[above,sloped,  font=\scriptsize](C,B){\SI{9}{\centi\meter}}  
	%	\tkzLabelPoints[below left,font=\scriptsize](A)
	%	\tkzLabelPoints[below right,font=\scriptsize](B)
	%	\tkzLabelPoints[above,font=\scriptsize](C)  
	\tkzDefPointBy[projection = onto A--B](C) \tkzGetPoint{H}
	\tkzDrawSegment[dashed](C,H)
	\tkzMarkRightAngle[size=0.2](B,H,C)
	%	\tkzLabelPoints[below left,font=\scriptsize](H)
\end{tikzpicture}   
} 

\noindent\parbox[position=t]{0.5\textwidth}{\centering
	%\begin{figure}[H]\centering
	\begin{tikzpicture}[scale=1.0] 
		\tkzSetUpPoint[size=3, fill=black]
		\tkzfctset{tan style/.style={->,>=latex, color=red, line width=1.2pt}}   
		\tkzSetUpLine[line width= 1.2pt, add=.5 and .5] 
		
		\tkzDefPoint(0,0){A} 
		\tkzDefShiftPoint[A](0:3){B}
		\tkzDefShiftPoint[A](50:3){C}
		% 			\tkzDefTriangle[two angles=90 and 35](A,B)\tkzGetPoint{C}  
		\tkzDrawPolygon[line width=1.2pt](A,B,C)
		%		\tkzLabelPoints[font=\scriptsize](A,B,C)
		% 			\tkzMarkRightAngle[](B,A,C)
		\tkzMarkAngle[mark=none,-, >=latex, size =0.5](B,A,C)
		% 			\tkzMarkAngle[mark=none,-, >=latex, size =0.5](C,B,A)
		%\tkzMarkAngle[mark=none,-, >=latex, size =0.5](A,C,B)
		\tkzLabelAngle[font=\footnotesize,pos=.9](B,A,C){\scriptsize$\ang{50}$}
		% 			\tkzLabelAngle[font=\footnotesize,pos=.9](C,B,A){$?$}
		%\tkzLabelAngle[font=\scriptsize,pos=.8](A,C,B){$\ang{55}$}
		
%		\tkzDrawSegment[dim={\scriptsize \SI{4}{\centi\meter},+15pt,}](A,C)
		\tkzDrawSegment[dim={\scriptsize \SI{5}{\centi\meter},-15pt,}](A,B)
		
		\tkzMarkSegments[mark=s||](A,B A,C)
		
		%		\tkzLabelSegment[above, font=\scriptsize,sloped](A,C){\SI{5}{\centi\meter}} 
		%		\tkzLabelSegment[below, font=\scriptsize](A,B){\SI{6}{\centi\meter}}
		
		% 			\tkzLabelSegment[above,sloped,  font=\scriptsize](C,B){\SI{9}{\centi\meter}}  
%		\tkzLabelPoints[below left,font=\scriptsize](A)
%		\tkzLabelPoints[below right,font=\scriptsize](B)
%		\tkzLabelPoints[above,font=\scriptsize](C)  
		\tkzDefPointBy[projection = onto A--B](C) \tkzGetPoint{H}
		\tkzDrawSegment[dashed](C,H)
		\tkzMarkRightAngle[size=0.2](B,H,C)
%		\tkzLabelPoints[below left,font=\scriptsize](H) 
		\begin{scope}[xshift=-2.5cm]
					\tkzSetUpPoint[size=3, fill=black]
			\tkzfctset{tan style/.style={->,>=latex, color=red, line width=1.2pt}}   
			\tkzSetUpLine[line width= 1.2pt, add=.5 and .5] 
			
			\tkzDefPoint(0,0){A} 
			\tkzDefShiftPoint[A](-180:3){B}
			\tkzDefShiftPoint[A](50:3){C}
			% 			\tkzDefTriangle[two angles=90 and 35](A,B)\tkzGetPoint{C}  
			\tkzDrawPolygon[line width=1.2pt](A,B,C)
			%		\tkzLabelPoints[font=\scriptsize](A,B,C)
			% 			\tkzMarkRightAngle[](B,A,C)
			\tkzMarkAngle[mark=none,-, >=latex, size =0.4](C,A,B)
			% 			\tkzMarkAngle[mark=none,-, >=latex, size =0.5](C,B,A)
			%\tkzMarkAngle[mark=none,-, >=latex, size =0.5](A,C,B)
			\tkzLabelAngle[font=\footnotesize,pos=.9](C,A,B){\scriptsize$\ang{130}$}
			\tkzLabelAngle[font=\footnotesize,pos=.9](H,A,C){\scriptsize$?$}
			% 			\tkzLabelAngle[font=\footnotesize,pos=.9](C,B,A){$?$}
			%\tkzLabelAngle[font=\scriptsize,pos=.8](A,C,B){$\ang{55}$}
			
%		\tkzDrawSegment[dim={\scriptsize \SI{4}{\centi\meter},+15pt,}](A,C)
%			\tkzDrawSegment[dim={\scriptsize \SI{4}{\centi\meter},-15pt,}](A,B) 
			\tkzMarkSegments[mark=s||](A,B A,C) 
			%		\tkzLabelSegment[above, font=\scriptsize,sloped](A,C){\SI{5}{\centi\meter}} 
			%		\tkzLabelSegment[below, font=\scriptsize](A,B){\SI{6}{\centi\meter}}
			
			% 			\tkzLabelSegment[above,sloped,  font=\scriptsize](C,B){\SI{9}{\centi\meter}}  
%			\tkzLabelPoints[below left,font=\scriptsize](A)
%			\tkzLabelPoints[below right,font=\scriptsize](B)
%			\tkzLabelPoints[above,font=\scriptsize](C)  
			\tkzDefPointBy[projection = onto A--B](C) \tkzGetPoint{H}
			\tkzMarkAngle[mark=none,-, >=latex, size =0.6](H,A,C)
			\tkzDrawSegment[dashed](C,H)
			\tkzDrawLine[dashed](A,H)
			\tkzMarkRightAngle[size=0.2](B,H,C)
%			\tkzLabelPoints[below left,font=\scriptsize](H) 
		\end{scope}
	\end{tikzpicture}   
}\hfill \parbox[position=t]{0.5\textwidth}{\centering
%\begin{figure}[H]\centering
\begin{tikzpicture}[scale=1.0] 
	\tkzSetUpPoint[size=3, fill=black]
	\tkzfctset{tan style/.style={->,>=latex, color=red, line width=1.2pt}}   
	\tkzSetUpLine[line width= 1.2pt, add=.5 and .5] 
	
	\tkzDefPoint(0,0){A} 
	\tkzDefShiftPoint[A](0:2.5){B}
	\tkzDefShiftPoint[A](-40:2.5){C}
	% 			\tkzDefTriangle[two angles=90 and 35](A,B)\tkzGetPoint{C}  
	\tkzDrawPolygon[line width=1.2pt](A,B,C)
	%		\tkzLabelPoints[font=\scriptsize](A,B,C)
	% 			\tkzMarkRightAngle[](B,A,C)
	\tkzMarkAngle[mark=none,-, >=latex, size =0.5](C,A,B)
	% 			\tkzMarkAngle[mark=none,-, >=latex, size =0.5](C,B,A)
	%\tkzMarkAngle[mark=none,-, >=latex, size =0.5](A,C,B)
	\tkzLabelAngle[font=\footnotesize,pos=.9](C,A,B){\scriptsize$\ang{40}$}
	% 			\tkzLabelAngle[font=\footnotesize,pos=.9](C,B,A){$?$}
	%\tkzLabelAngle[font=\scriptsize,pos=.8](A,C,B){$\ang{55}$}
	
	%		\tkzDrawSegment[dim={\scriptsize \SI{4}{\centi\meter},+15pt,}](A,C)
	\tkzDrawSegment[dim={\scriptsize \SI{4}{\centi\meter},15pt,}](A,B)
	
	\tkzMarkSegments[mark=s|](A,B A,C)
	
	%		\tkzLabelSegment[above, font=\scriptsize,sloped](A,C){\SI{5}{\centi\meter}} 
	%		\tkzLabelSegment[below, font=\scriptsize](A,B){\SI{6}{\centi\meter}}
	
	% 			\tkzLabelSegment[above,sloped,  font=\scriptsize](C,B){\SI{9}{\centi\meter}}  
	%		\tkzLabelPoints[below left,font=\scriptsize](A)
	%		\tkzLabelPoints[below right,font=\scriptsize](B)
	%		\tkzLabelPoints[above,font=\scriptsize](C)  
	\tkzDefPointBy[projection = onto A--B](C) \tkzGetPoint{H}
	\tkzDrawSegment[dashed](C,H)
	\tkzMarkRightAngle[size=0.2](B,H,C)
	%		\tkzLabelPoints[below left,font=\scriptsize](H) 
	\begin{scope}[xshift=-2.5cm]
		\tkzSetUpPoint[size=3, fill=black]
		\tkzfctset{tan style/.style={->,>=latex, color=red, line width=1.2pt}}   
		\tkzSetUpLine[line width= 1.2pt, add=.5 and .5] 
		
		\tkzDefPoint(0,0){A} 
		\tkzDefShiftPoint[A](-180:2.5){B}
		\tkzDefShiftPoint[A](-40:2.5){C}
		% 			\tkzDefTriangle[two angles=90 and 35](A,B)\tkzGetPoint{C}  
		\tkzDrawPolygon[line width=1.2pt](A,B,C)
		%		\tkzLabelPoints[font=\scriptsize](A,B,C)
		% 			\tkzMarkRightAngle[](B,A,C)
		\tkzMarkAngle[mark=none,-, >=latex, size =0.4](B,A,C)
		% 			\tkzMarkAngle[mark=none,-, >=latex, size =0.5](C,B,A)
		%\tkzMarkAngle[mark=none,-, >=latex, size =0.5](A,C,B)
		\tkzLabelAngle[font=\footnotesize,pos=.7](B,A,C){\scriptsize$\ang{140}$}
		% 			\tkzLabelAngle[font=\footnotesize,pos=.9](C,B,A){$?$}
		%\tkzLabelAngle[font=\scriptsize,pos=.8](A,C,B){$\ang{55}$}
		
		%		\tkzDrawSegment[dim={\scriptsize \SI{4}{\centi\meter},+15pt,}](A,C)
		%			\tkzDrawSegment[dim={\scriptsize \SI{4}{\centi\meter},-15pt,}](A,B)
		
		\tkzMarkSegments[mark=s|](A,B A,C)
		
		%		\tkzLabelSegment[above, font=\scriptsize,sloped](A,C){\SI{5}{\centi\meter}} 
		%		\tkzLabelSegment[below, font=\scriptsize](A,B){\SI{6}{\centi\meter}}
		
		% 			\tkzLabelSegment[above,sloped,  font=\scriptsize](C,B){\SI{9}{\centi\meter}}  
		%			\tkzLabelPoints[below left,font=\scriptsize](A)
		%			\tkzLabelPoints[below right,font=\scriptsize](B)
		%			\tkzLabelPoints[above,font=\scriptsize](C)  
		\tkzDefPointBy[projection = onto A--B](C) \tkzGetPoint{H}
		\tkzDrawSegment[dashed](C,H)
		\tkzDrawLine[dashed](A,H)
		\tkzMarkRightAngle[size=0.2](B,H,C)
		%			\tkzLabelPoints[below left,font=\scriptsize](H) 
	\end{scope}
\end{tikzpicture}   
}    
\end{exercise}
 \begin{remark}
 	On peut calculer des rapports trigonométriques d'angles obtus (>\ang{90}). Comparer $\sin(\ang{130})$,$\sin(\ang{50})$. 
 \end{remark}
\newpage 

\begin{tBox}[frametitle=Formule de l'aire]

\noindent\parbox[position=t]{0.35\textwidth}{\centering
	%\begin{figure}[H]\centering
	\begin{tikzpicture}[scale=1.0] 
		\tkzSetUpPoint[size=3, fill=black]
		\tkzfctset{tan style/.style={->,>=latex, color=red, line width=1.2pt}}   
		\tkzSetUpLine[line width= 1.2pt, add=.5 and .5] 
		
		\tkzDefPoint(0,0){C} 
		\tkzDefShiftPoint[C](-10:4.2){A}
		\tkzDefShiftPoint[C](45:3){B}
		% 			\tkzDefTriangle[two angles=90 and 35](A,B)\tkzGetPoint{C}  
		\tkzDrawPolygon[line width=1.2pt](A,B,C)
		%		\tkzLabelPoints[font=\scriptsize](A,B,C)
		% 			\tkzMarkRightAngle[](B,A,C)
		\tkzMarkAngle[mark=none,-, >=latex, size =0.5](A,C,B)
		% 			\tkzMarkAngle[mark=none,-, >=latex, size =0.5](C,B,A)
		%\tkzMarkAngle[mark=none,-, >=latex, size =0.5](A,C,B)
		\tkzLabelAngle[font=\footnotesize,pos=.9](A,C,B){$\widehat{C}$}
		% 			\tkzLabelAngle[font=\footnotesize,pos=.9](C,B,A){$?$}
		%\tkzLabelAngle[font=\scriptsize,pos=.8](A,C,B){$\ang{55}$}
		
		\tkzDrawSegment[dim={\scriptsize $b$,-15pt,}](C,A)
		\tkzDrawSegment[dim={\scriptsize $a$,+15pt,}](C,B)
		
		
		%		\tkzLabelSegment[above, font=\scriptsize,sloped](A,C){\SI{5}{\centi\meter}} 
		%		\tkzLabelSegment[below, font=\scriptsize](A,B){\SI{6}{\centi\meter}}
		
		% 			\tkzLabelSegment[above,sloped,  font=\scriptsize](C,B){\SI{9}{\centi\meter}}  
		\tkzLabelPoints[below left,font=\scriptsize](C)
		\tkzLabelPoints[below right,font=\scriptsize](A)
		\tkzLabelPoints[above,font=\scriptsize](B)  
		\tkzDefPointBy[projection = onto C--B](A) \tkzGetPoint{H}
		\tkzDrawSegment[dashed](C,H)
		\tkzMarkRightAngle[size=0.2](B,H,C)
%		\tkzLabelPoints[below left,font=\scriptsize](H)
	\end{tikzpicture}   
}\hfill 
\parbox[position=t]{0.6\textwidth}{
	La formule est valable pour des angles aigus, mais aussi des angles obtus !
	\[
		\mathscr{Aire}=\frac{1}{2}ab\sin(\widehat{C})
	\]
	On a aussi l'égalité : $\sin(\widehat{C})=\sin(\ang{180}-\widehat{C})$
}
\end{tBox}

\begin{example}[je fais : loi des sinus]\hfill 
	
\noindent\parbox[position=t]{0.25\textwidth}{\centering
	%\begin{figure}[H]\centering
	\begin{tikzpicture}[scale=1.0] 
		\tkzSetUpPoint[size=3, fill=black]
		\tkzfctset{tan style/.style={->,>=latex, color=red, line width=1.2pt}}   
		\tkzSetUpLine[line width= 1.2pt, add=.5 and .5] 
		
		\tkzDefPoint(0,0){C} 
		\tkzDefShiftPoint[C](-10:4.2){A}
		\tkzDefShiftPoint[C](45:3){B}
		% 			\tkzDefTriangle[two angles=90 and 35](A,B)\tkzGetPoint{C}  
		\tkzDrawPolygon[line width=1.2pt](A,B,C)
		%		\tkzLabelPoints[font=\scriptsize](A,B,C)
		% 			\tkzMarkRightAngle[](B,A,C)
		\tkzMarkAngle[mark=none,-, >=latex, size =0.5](A,C,B)
		\tkzMarkAngle[mark=none,-, >=latex, size =0.5](B,A,C)
		\tkzMarkAngle[mark=none,-, >=latex, size =0.5](C,B,A)
		% 			\tkzMarkAngle[mark=none,-, >=latex, size =0.5](C,B,A)
		%\tkzMarkAngle[mark=none,-, >=latex, size =0.5](A,C,B)
		\tkzLabelAngle[font=\footnotesize,pos=.9](A,C,B){\scriptsize$\widehat{C}$}
		\tkzLabelAngle[font=\footnotesize,pos=.9](B,A,C){\scriptsize$\widehat{A}$}
		\tkzLabelAngle[font=\footnotesize,pos=.9](C,B,A){\scriptsize$\widehat{B}$}
		% 			\tkzLabelAngle[font=\footnotesize,pos=.9](C,B,A){$?$}
		%\tkzLabelAngle[font=\scriptsize,pos=.8](A,C,B){$\ang{55}$}
		
		\tkzDrawSegment[dim={\scriptsize $b$,-15pt,}](C,A)
		\tkzDrawSegment[dim={\scriptsize $a$,+15pt,}](C,B)
		\tkzDrawSegment[dim={\scriptsize $c$,-15pt,}](A,B)
		
		
		%		\tkzLabelSegment[above, font=\scriptsize,sloped](A,C){\SI{5}{\centi\meter}} 
		%		\tkzLabelSegment[below, font=\scriptsize](A,B){\SI{6}{\centi\meter}}
		
		% 			\tkzLabelSegment[above,sloped,  font=\scriptsize](C,B){\SI{9}{\centi\meter}}  
		\tkzLabelPoints[below left,font=\small](C)
		\tkzLabelPoints[below right,font=\small](A)
		\tkzLabelPoints[above,font=\small](B)    
	\end{tikzpicture}   
}  
\parbox[position=t]{0.55\textwidth}{ 
$\begin{WithArrows}
2\mathscr{Aire}&= ab\sin(\widehat{C}) =  \\
&\\
&\\
&
\end{WithArrows}$   
} 

\noindent\parbox[position=t]{0.25\textwidth}{\centering 
	\begin{tikzpicture}[scale=1.0] 
		\tkzSetUpPoint[size=3, fill=black]
		\tkzfctset{tan style/.style={->,>=latex, color=red, line width=1.2pt}}   
		\tkzSetUpLine[line width= 1.2pt, add=.5 and .5]  
		\tkzDefPoint(0,0){C} 
		\tkzDefShiftPoint[C](-10:2.1){A}
		\tkzDefShiftPoint[C](45:1.5){B}  
		\tkzDrawPolygon[line width=1.2pt](A,B,C) 
		\tkzDrawSegment[dim={\scriptsize $\sin(\widehat{B})$,-20pt,}](C,A)
		\tkzDrawSegment[dim={\scriptsize $\sin(\widehat{A})$,+20pt,}](C,B)
		\tkzDrawSegment[dim={\scriptsize $\sin(\widehat{C})$,-20pt,}](A,B)  
	\end{tikzpicture}   
}\hfill 
\end{example}
\begin{exercise}\label{chapitre05:exercice:formuledesinus01}\hfill 
	
On considère la figure ci-dessous. Cocher les cases pour les formules vraies.

\noindent\parbox[position=t]{0.25\textwidth}{\centering 
	\begin{tikzpicture}[scale=1.0] 
		\tkzSetUpPoint[size=3, fill=black]
		\tkzfctset{tan style/.style={->,>=latex, color=red, line width=1.2pt}}   
		\tkzSetUpLine[line width= 1.2pt, add=.5 and .5]  
		\tkzDefPoint(0,0){C} 
		\tkzDefShiftPoint[C](-20:3.0){A}
		\tkzDefShiftPoint[C](35:2.0){B}  
		\tkzDrawPolygon[line width=1.2pt](A,B,C) 
		\tkzDrawSegment[dim={\scriptsize $c$,-20pt,}](C,A)
		\tkzDrawSegment[dim={\scriptsize $a$,+20pt,}](C,B)
		\tkzDrawSegment[dim={\scriptsize $b$,-20pt,}](A,B) 
		\tkzLabelAngle[font=\footnotesize,pos=.8](A,C,B){\scriptsize$\widehat{d}$}
		\tkzLabelAngle[font=\footnotesize,pos=.8](B,A,C){\scriptsize$\widehat{f}$}
		\tkzLabelAngle[font=\footnotesize,pos=.8](C,B,A){\scriptsize$\widehat{e}$}  
		\tkzMarkAngle[mark=none,-, >=latex, size =0.5](A,C,B)
		\tkzMarkAngle[mark=none,-, >=latex, size =0.5](B,A,C)
		\tkzMarkAngle[mark=none,-, >=latex, size =0.5](C,B,A)   
	\end{tikzpicture}   
}\hfill 
\parbox[position=t]{0.7\textwidth}{
\QCM[VF,Alterne,Stretch=1.7, Largeur=0.75cm]{%
	{$\frac{\sin(d)}{b}=\frac{\sin(e)}{c}$}&1,%
	{$\frac{a}{\sin(f)}=\frac{b}{\sin(c)}$}&2,%
	{$\frac{c}{\sin(e)}=\frac{b}{\sin(d)}$}&1,%
} 
}
\end{exercise}
\begin{exercise}\label{chapitre05:exercice:formuledesinus02}\hfill 
	
En utilisant la loi des sinus, trouver les valeurs de $x$, $y$ et $z$ pour chacune des figures suivantes.

\noindent\parbox[position=t]{0.30\textwidth}{\centering 
	\begin{tikzpicture}[scale=1.0] 
		\tkzSetUpPoint[size=3, fill=black]
		\tkzfctset{tan style/.style={->,>=latex, color=red, line width=1.2pt}}   
		\tkzSetUpLine[line width= 1.2pt, add=.5 and .5]  
		\tkzDefPoint(0,0){C} 
		\tkzDefShiftPoint[C](0:3.3){A}
		\tkzDefShiftPoint[C](52:2.33){B}  
		
		
		\tkzDrawPolygon[line width=1.2pt](A,B,C) 
%		\tkzDrawSegment[dim={\scriptsize $c$,-20pt,}](C,A)
		\tkzDrawSegment[dim={\small \SI{7}{\centi\meter},+20pt,}](C,B)
		\tkzDrawSegment[dim={\small $x$,-20pt,}](A,B) 
		\tkzLabelAngle[font=\footnotesize,pos=.8](A,C,B){\scriptsize$\ang{52}$}
		\tkzLabelAngle[font=\footnotesize,pos=.8](B,A,C){\scriptsize$\ang{45}$}
%		\tkzLabelAngle[font=\footnotesize,pos=.8](C,B,A){\scriptsize$\widehat{e}$}  
		\tkzMarkAngle[mark=none,-, >=latex, size =0.5](A,C,B)
		\tkzMarkAngle[mark=none,-, >=latex, size =0.5](B,A,C)
%		\tkzMarkAngle[mark=none,-, >=latex, size =0.5](C,B,A)   
	\end{tikzpicture}   
}\hfill 
\noindent\parbox[position=t]{0.3\textwidth}{\centering 
	\begin{tikzpicture}[scale=1.0] 
		\tkzSetUpPoint[size=3, fill=black]
		\tkzfctset{tan style/.style={->,>=latex, color=red, line width=1.2pt}}   
		\tkzSetUpLine[line width= 1.2pt, add=.5 and .5]  
		\tkzDefPoint(0,0){C} 
		\tkzDefShiftPoint[C](-20:4.2){A}
		\tkzDefShiftPoint[C](36:3.0){B}  
		\tkzDrawPolygon[line width=1.2pt](A,B,C) 
%		\tkzDrawSegment[dim={\scriptsize $c$,-20pt,}](C,A)
		\tkzDrawSegment[dim={\small \SI{6}{\centi\meter},+20pt,}](C,B)
		\tkzDrawSegment[dim={\small \SI{7}{\centi\meter},-20pt,}](A,B) 
		\tkzLabelAngle[font=\footnotesize,pos=.8](A,C,B){\scriptsize$y^\circ$}
		\tkzLabelAngle[font=\footnotesize,pos=.8](B,A,C){\scriptsize$\ang{45}$}
%		\tkzLabelAngle[font=\footnotesize,pos=.8](C,B,A){\scriptsize$\widehat{e}$}  
		\tkzMarkAngle[mark=none,-, >=latex, size =0.5](A,C,B)
		\tkzMarkAngle[mark=none,-, >=latex, size =0.5](B,A,C)
%		\tkzMarkAngle[mark=none,-, >=latex, size =0.5](C,B,A)   
	\end{tikzpicture}   
}\hfill
\noindent\parbox[position=t]{0.3\textwidth}{\centering 
	\begin{tikzpicture}[scale=1.0] 
		\tkzSetUpPoint[size=3, fill=black]
		\tkzfctset{tan style/.style={->,>=latex, color=red, line width=1.2pt}}   
		\tkzSetUpLine[line width= 1.2pt, add=.5 and .5]  
		\tkzDefPoint(0,0){C} 
		\tkzDefShiftPoint[C](-20:3.0){A}
		\tkzDefShiftPoint[C](35:2.0){B}  
		\tkzDrawPolygon[line width=1.2pt](A,B,C) 
		\tkzDrawSegment[dim={\small \SI{8}{\centi\meter},-20pt,}](C,A)
%		\tkzDrawSegment[dim={\scriptsize $a$,+20pt,}](C,B)
		\tkzDrawSegment[dim={\small $z$,-20pt,}](A,B) 
		\tkzLabelAngle[font=\footnotesize,pos=.8](A,C,B){\scriptsize$\ang{52}$}
		\tkzLabelAngle[font=\footnotesize,pos=.8](B,A,C){\scriptsize$\ang{45}$}
%		\tkzLabelAngle[font=\footnotesize,pos=.8](C,B,A){\scriptsize$\widehat{e}$}  
		\tkzMarkAngle[mark=none,-, >=latex, size =0.5](A,C,B)
		\tkzMarkAngle[mark=none,-, >=latex, size =0.5](B,A,C)
%		\tkzMarkAngle[mark=none,-, >=latex, size =0.5](C,B,A)   
	\end{tikzpicture}   
}
\end{exercise}
%----------------------------------------------------------------------------------------
%	 
%----------------------------------------------------------------------------------------
\newpage
\begin{proof}[solution de l'exercice \ref{chapter05:exercices:trigonometrie:01}] \hfill 	
\begin{multicols}{3}
	\begin{enumerate}[label=\Alph*),leftmargin=*]
		\item $\cos(x)=\frac{7}{9}$.\\
		$x= \arccos(\tfrac{7}{9})\approx\ang{39}$.
		\item $\sin(x)=\frac{4}{5}$.\\
		$x= \arcsin(\tfrac{4}{5})\approx \ang{53}$
		\item $x=5\cos(\ang{40})\approx \numprint{3.83}$.
		\item $x=\frac{5}{\cos(\ang{40})}\approx\numprint{6.53}$.
		\item $x=\frac{3}{\sin(\ang{40})}\approx \numprint{4.67}$.
		\item $\tan(x)=\frac{4}{3}$.\\
		$x=\arctan(\tfrac{4}{3})\approx\ang{53.1}$.
		\item $\cos(x)=\frac{5}{7}$\\
			$x=\arccos(\tfrac{5}{7})=\ang{44.4}$.
		\item $\sin(x)=\frac{5}{10}$, $x=\ang{30}$.
	\end{enumerate}
\end{multicols} 
\end{proof}

\begin{proof}[solution de l'exercice \ref{chapitre05:exercice:calculdaires01}] \hfill 	
	
$A\approx\SI{13.6}{\centi\meter^2}$; $B\approx\SI{18.5}{\centi\meter^2}$; $C\approx\SI{29.1}{\centi\meter^2}$; $D\approx\SI{9.58}{\centi\meter^2}$; $E\approx\SI{5.14}{\centi\meter^2}$; 
\end{proof}

\begin{proof}[solution de l'exercice \ref{chapitre05:exercice:formuledesinus01}] \hfill 	

\QCM[VF,Alterne,Stretch=1.7, Solution, Largeur=0.75cm]{%
	{$\frac{\sin(d)}{b}=\frac{\sin(e)}{c}$}&1,%
	{$\frac{a}{\sin(f)}=\frac{b}{\sin(c)}$}&2,%
	{$\frac{c}{\sin(e)}=\frac{b}{\sin(d)}$}&1,%
} 
\end{proof}
\begin{proof}[solution de l'exercice \ref{chapitre05:exercice:formuledesinus02}] \hfill 	
	
\noindent $x=\SI{6.3}{\centi\meter} $, $y\approx \ang{55.6}$ et $z\approx \SI{6.35}{\centi\meter}$
\end{proof}


%----------------------------------------------------------------------------------------
%	EXERCICES
%----------------------------------------------------------------------------------------
 
%----------------------------------------------------------------------------------------
%	
%----------------------------------------------------------------------------------------
\newpage 
\subsection{La loi des cosinus}

\begin{exercise}[un classique de 3\ieme]\hfill 
	
	\noindent\parbox[position=t]{0.35\textwidth}{\centering
		%\begin{figure}[H]\centering
		\begin{tikzpicture}[scale=1.0] 
			\tkzSetUpPoint[size=3, fill=black]
			\tkzfctset{tan style/.style={->,>=latex, color=red, line width=1.2pt}}   
			\tkzSetUpLine[line width= 1.2pt, add=.5 and .5] 
			
			\tkzDefPoint(0,0){A} 
			\tkzDefShiftPoint[A](0:5){B}
			\tkzDefShiftPoint[A](38:3.0){C}
			% 			\tkzDefTriangle[two angles=90 and 35](A,B)\tkzGetPoint{C}  
			\tkzDrawPolygon[line width=1.2pt](A,B,C)
			%		\tkzLabelPoints[font=\scriptsize](A,B,C)
			% 			\tkzMarkRightAngle[](B,A,C)
			\tkzMarkAngle[mark=none,-, >=latex, size =0.5](B,A,C)
			% 			\tkzMarkAngle[mark=none,-, >=latex, size =0.5](C,B,A)
			%\tkzMarkAngle[mark=none,-, >=latex, size =0.5](A,C,B)
			\tkzLabelAngle[font=\footnotesize,pos=.9](B,A,C){$\ang{38}$}
			% 			\tkzLabelAngle[font=\footnotesize,pos=.9](C,B,A){$?$}
			%\tkzLabelAngle[font=\scriptsize,pos=.8](A,C,B){$\ang{55}$}
			\tkzLabelSegment[above, font=\scriptsize,sloped](A,C){\SI{6}{\centi\meter}} 
			\tkzLabelSegment[below, font=\scriptsize](A,B){\SI{10}{\centi\meter}}
			\tkzLabelSegment[above,sloped,  font=\scriptsize](C,B){$x$}  
			\tkzLabelPoints[below left,font=\scriptsize](A)
			\tkzLabelPoints[below right,font=\scriptsize](B)
			\tkzLabelPoints[above,font=\scriptsize](C)  
			\tkzDefPointBy[projection = onto A--B](C) \tkzGetPoint{H}
			\tkzDrawSegment[dashed](C,H)
			\tkzMarkRightAngle[size=0.2](B,H,C)
			\tkzLabelPoints[above left,font=\scriptsize](H)
		\end{tikzpicture}   
	}  \hfill 
	\parbox[position=t]{0.6\textwidth}{
		Dans la figure ci-contre, $H$ est le pied de la hauteur issue de $C$.
		\begin{enumerate}[label=\alph*),leftmargin=*]
			\item Calculer les longueurs $CH$ et $HB$.
			\item En déduire $HB$, puis une valeur approchée de $x$ au dixième de \si{\centi\meter}.
		\end{enumerate}	
	}
\end{exercise}	
\begin{example}[je fais, loi des cosinus]\hfill 
	
	
	\noindent\parbox[position=t]{0.4\textwidth}{\centering 
		\begin{tikzpicture}[scale=1.0] 
			\tkzSetUpPoint[size=3, fill=black]
			\tkzfctset{tan style/.style={->,>=latex, color=red, line width=1.2pt}}   
			\tkzSetUpLine[line width= 1.2pt, add=.5 and .5]  
			\tkzDefPoint(0,0){C} 
			\tkzDefShiftPoint[C](0:6.10){A}
			\tkzDefShiftPoint[C](52:2.70){B}  
			\tkzDrawPolygon[line width=1.2pt](A,B,C) 
			\tkzDrawSegment[dim={$b$,-10pt,}](C,A)
			\tkzDrawSegment[dim={$a$,+10pt,}](C,B)
			\tkzDrawSegment[dim={$c=?$,-10pt,}](A,B) 
			\tkzLabelAngle[font=\small,pos=.8](A,C,B){$\widehat{C}$}  
			\tkzMarkAngle[mark=none,-, >=latex, size =0.5](A,C,B)  
			\tkzDefPointBy[projection = onto C--A](B) \tkzGetPoint{H}
			\tkzDrawSegment[dashed](B,H)
			\tkzMarkRightAngle[size=0.2](A,H,B)  
		\end{tikzpicture}   
	} 
	\vfill 
\end{example}
\begin{exercise}\hfill 
	
	\noindent\parbox[position=t]{0.25\textwidth}{\centering 
		\begin{tikzpicture}[scale=1.0] 
			\tkzSetUpPoint[size=3, fill=black]
			\tkzfctset{tan style/.style={->,>=latex, color=red, line width=1.2pt}}   
			\tkzSetUpLine[line width= 1.2pt, add=.5 and .5]  
			\tkzDefPoint(0,0){C} 
			\tkzDefShiftPoint[C](-10:3.50){A}
			\tkzDefShiftPoint[C](45:2.10){B}  
			\tkzDrawPolygon[line width=1.2pt](A,B,C) 
			\tkzDrawSegment[dim={$\SI{6}{\centi\meter}$,-10pt,}](C,A)
			\tkzDrawSegment[dim={$a$,+20pt,}](C,B)
			\tkzDrawSegment[dim={$\SI{5}{\centi\meter}$,-10pt,}](A,B) 
			%		\tkzLabelAngle[font=\small,pos=.8](A,C,B){$\ang{40}$} 
			%		\tkzLabelAngle[font=\small,pos=.8](C,A,B){$\ang{40}$}  
			\tkzLabelAngle[font=\small,pos=.9](B,A,C){$\ang{40}$}  
			\tkzMarkAngle[mark=none,-, >=latex, size =0.5](B,A,C)    
		\end{tikzpicture}   
	} \hfill 
	\parbox[position=t]{0.7\textwidth}{  
		$\begin{WithArrows}
			a^2&=\rule{0pt}{1.8em} 5^2+6^2-2\times 5\times 6 \times \cos(\ang{40})  \\
			a& =\rule{0pt}{1.8em}  
		\end{WithArrows}$  
	}
	
	\noindent\parbox[position=t]{0.25\textwidth}{\centering 
		\begin{tikzpicture}[scale=1.0] 
			\tkzSetUpPoint[size=3, fill=black]
			\tkzfctset{tan style/.style={->,>=latex, color=red, line width=1.2pt}}   
			\tkzSetUpLine[line width= 1.2pt, add=.5 and .5]  
			\tkzDefPoint(0,0){C} 
			\tkzDefShiftPoint[C](5:4.10){A}
			\tkzDefShiftPoint[C](-38:2.70){B}  
			\tkzDrawPolygon[line width=1.2pt](A,B,C) 
			\tkzDrawSegment[dim={$b$,+10pt,}](C,A)
			\tkzDrawSegment[dim={$\SI{6}{\centi\meter}$,-10pt,}](C,B)
			\tkzDrawSegment[dim={$\SI{5}{\centi\meter}$,+10pt,}](A,B)   
			\tkzLabelAngle[font=\small,pos=.9](A,B,C){$\ang{85}$}  
			\tkzMarkAngle[mark=none,-, >=latex, size =0.5](A,B,C)    
		\end{tikzpicture}   
	} \hfill 
	\parbox[position=t]{0.7\textwidth}{  
		$\begin{WithArrows}
			b^2&=\\
			b& =\rule{0pt}{1.8em}  
		\end{WithArrows}$  
	}
	
	\noindent\parbox[position=t]{0.25\textwidth}{\centering 
		\begin{tikzpicture}[scale=1.0] 
			\tkzSetUpPoint[size=3, fill=black]
			\tkzfctset{tan style/.style={->,>=latex, color=red, line width=1.2pt}}   
			\tkzSetUpLine[line width= 1.2pt, add=.5 and .5]  
			\tkzDefPoint(0,0){C} 
			\tkzDefShiftPoint[C](-17.5:3){A}
			\tkzDefShiftPoint[C](-90:2){B}  
			\tkzDrawPolygon[line width=1.2pt](A,B,C) 
			\tkzDrawSegment[dim={$\SI{4.5}{\centi\meter}$,+10pt,}](C,A)
			\tkzDrawSegment[dim={$c$,-10pt,}](C,B)
			%		\tkzDrawSegment[dim={$ $,+10pt,}](A,B)   
			\tkzLabelAngle[font=\small,pos=.9](C,A,B){$\ang{35}$}  
			\tkzMarkAngle[mark=none,-, >=latex, size =0.5](C,A,B)    
			\tkzMarkSegments[mark=s||](A,B A,C)
		\end{tikzpicture}   
	} \hfill 
	\parbox[position=t]{0.7\textwidth}{  
		$\begin{WithArrows}
			c^2&=\\
			c& =\rule{0pt}{1.8em}  
		\end{WithArrows}$  
	}
	
	\noindent\parbox[position=t]{0.25\textwidth}{\centering 
		\begin{tikzpicture}[scale=1.0] 
			\tkzSetUpPoint[size=3, fill=black]
			\tkzfctset{tan style/.style={->,>=latex, color=red, line width=1.2pt}}   
			\tkzSetUpLine[line width= 1.2pt, add=.5 and .5]  
			\tkzDefPoint(0,0){C} 
			\tkzDefShiftPoint[C](-17.5:3){A}
			\tkzDefShiftPoint[C](-90:2){B}  
			\tkzDrawPolygon[line width=1.2pt](A,B,C) 
			\tkzDrawSegment[dim={$\SI{4.5}{\centi\meter}$,+10pt,}](C,A)
			\tkzDrawSegment[dim={$c$,-10pt,}](C,B)
			%		\tkzDrawSegment[dim={$ $,+10pt,}](A,B)   
			\tkzLabelAngle[font=\small,pos=.9](C,A,B){$\ang{35}$}  
			\tkzMarkAngle[mark=none,-, >=latex, size =0.5](C,A,B)    
			\tkzMarkSegments[mark=s||](A,B A,C)
		\end{tikzpicture}   
	} \hfill   
	
	
	%----------------------------------------------------------------------------------------
	%	
	%----------------------------------------------------------------------------------------
	\newpage 
	
	\noindent\parbox[position=t]{0.25\textwidth}{\centering 
		\begin{tikzpicture}[scale=1.0] 
			\tkzSetUpPoint[size=3, fill=black]
			\tkzfctset{tan style/.style={->,>=latex, color=red, line width=1.2pt}}   
			\tkzSetUpLine[line width= 1.2pt, add=.5 and .5]  
			\tkzDefPoint(0,0){C} 
			\tkzDefShiftPoint[C](-10:3.5){A}
			\tkzDefShiftPoint[C](42:2.5){B}  
			\tkzDrawPolygon[line width=1.2pt](A,B,C) 
			\tkzDrawSegment[dim={$\SI{5.2}{\centi\meter}$,-10pt,}](C,A)
			\tkzDrawSegment[dim={$\SI{3}{\centi\meter}$,+10pt,}](C,B)
			\tkzDrawSegment[dim={$\SI{8}{\centi\meter}$,-10pt,}](A,B) 
			\tkzLabelAngle[font=\small,pos=.8](A,C,B){$A$}  
			\tkzMarkAngle[mark=none,-, >=latex, size =0.5](A,C,B)  
			%		\tkzDefPointBy[projection = onto C--A](B) \tkzGetPoint{H}
			%		\tkzDrawSegment[dashed](B,H)
			%		\tkzMarkRightAngle[size=0.2](A,H,B)  
		\end{tikzpicture}   
	} \hfill 
	\parbox[position=t]{0.7\textwidth}{  
		$\begin{WithArrows}
			\big( \qquad \big)^2&=\big( \qquad \big)^2+\big( \qquad \big)^2-2\times \qquad \times \qquad \times \cos(A) \\
			&= \rule{0pt}{1.8em} \\
			\cos(A)& =\rule{0pt}{1.8em} \\
			A& =\rule{0pt}{1.8em}  
		\end{WithArrows}$  
	}
	\bigskip
	
	\noindent  \parbox[position=t]{0.30\textwidth}{\centering 
		\begin{tikzpicture}[scale=1.0] 
			\tkzSetUpPoint[size=3, fill=black]
			\tkzfctset{tan style/.style={->,>=latex, color=red, line width=1.2pt}}   
			\tkzSetUpLine[line width= 1.2pt, add=.5 and .5]  
			\tkzDefPoint(0,0){C} 
			\tkzDefShiftPoint[C](0:3.5){A}
			\tkzDefShiftPoint[C](115:2.5){B}  
			\tkzDrawPolygon[line width=1.2pt](A,B,C) 
			\tkzDrawSegment[dim={$\SI{4.3}{\centi\meter}$,-10pt,}](C,A)
			\tkzDrawSegment[dim={$\SI{2.6}{\centi\meter}$,+10pt,}](C,B)
			\tkzDrawSegment[dim={$\SI{6.1}{\centi\meter}$,-10pt,}](A,B) 
			%		\tkzLabelAngle[font=\small,pos=.8](A,C,B){$A$}  
			\tkzLabelAngle[font=\small,pos=.9](B,A,C){$B$} 
			\tkzMarkAngle[mark=none,-, >=latex, size =0.6](B,A,C)  
			%		\tkzDefPointBy[projection = onto C--A](B) \tkzGetPoint{H}
			%		\tkzDrawSegment[dashed](B,H)
			%		\tkzMarkRightAngle[size=0.2](A,H,B)  
		\end{tikzpicture}   
	} \bigskip 
	
	
	\noindent \parbox[position=t]{0.30\textwidth}{\centering 
		\begin{tikzpicture}[scale=1.0] 
			\tkzSetUpPoint[size=3, fill=black]
			\tkzfctset{tan style/.style={->,>=latex, color=red, line width=1.2pt}}   
			\tkzSetUpLine[line width= 1.2pt, add=.5 and .5]  
			\tkzDefPoint(0,0){C} 
			\tkzDefShiftPoint[C](-107.5:3){A}
			\tkzDefShiftPoint[C](-180:2){B}  
			\tkzDrawPolygon[line width=1.2pt](A,B,C) 
			\tkzDrawSegment[dim={$\SI{1.8}{\centi\meter}$,+25pt,}](C,A)
			\tkzDrawSegment[dim={$\SI{1.3}{\centi\meter}$,-10pt,}](C,B)
			%		\tkzDrawSegment[dim={$ $,+10pt,}](A,B) 
			\tkzMarkAngle[mark=none,-, >=latex, size =0.6](C,A,B)    
			\tkzLabelAngle[font=\small,pos=.9](C,A,B){$C$}    
			\tkzMarkSegments[mark=s||](A,B A,C)
		\end{tikzpicture}   
	}  
\end{exercise}
\begin{exercise}\hfill 
	
	\noindent\parbox[position=t]{0.25\textwidth}{\centering
		%\begin{figure}[H]\centering
		\begin{tikzpicture}[scale=1.0] 
			\tkzSetUpPoint[size=3, fill=black]
			\tkzfctset{tan style/.style={->,>=latex, color=red, line width=1.2pt}}   
			\tkzSetUpLine[line width= 1.2pt, add=.5 and .5] 
			
			\tkzDefPoint(0,0){P} 
			\tkzDefShiftPoint[P](-5:3.6){Q}
			\tkzDefShiftPoint[P](71:2.8){R} 
			\tkzDrawPolygon[line width=1.2pt](P,Q,R) 
			\tkzMarkAngle[mark=none,-, >=latex, size =0.5](Q,P,R) 
			\tkzLabelAngle[font=\footnotesize,pos=.9](Q,P,R){$\ang{76}$} 
			\tkzLabelSegment[above, font=\scriptsize,sloped](P,R){\SI{7}{\centi\meter}} 
			\tkzLabelSegment[below, font=\scriptsize](P,Q){\SI{9}{\centi\meter}}
			
			\tkzLabelPoints[below left,font=\scriptsize](P)
			\tkzLabelPoints[below right,font=\scriptsize](Q)
			\tkzLabelPoints[above,font=\scriptsize](R)    
		\end{tikzpicture}   
	}\hfill 
	\parbox[position=t]{0.7\textwidth}{
		Dans le triangle $PQR$, $\widehat{QPR}= \ang{76}$, $PQ=\SI{9}{\centi\meter}$ et $PR=\SI{7}{\centi\meter}$.  
		\begin{enumerate}[label=\alph*), leftmargin=*]
			\item Calculer $QR$ au millimètre près.
			\item À l'aide de la loi des sinus, calculer les valeurs approchées à $10^{-1}$ près des angles $\widehat{PQR}$ et $\widehat{PRQ}$. 
		\end{enumerate}
		
		
	} 
\end{exercise}

\cleardoublepage
 

%----------------------------------------------------------------------------------------
%	 
%----------------------------------------------------------------------------------------
\end{document}